%!TEX root = main.tex

\section{Asymptotic Normality: Complete Proof}

\subsection{Setup and Notation}

\paragraph{Estimand.}
Let $\mu$ be a distribution over probability measures $\mathcal{Q}$ on the sample space $(\mathcal{O}, \mathcal{A})$. For a given distribution $\mathcal{Q}$, define treatment effects:
\begin{align}
\Delta_S(\mathcal{Q}) &= \mathbb{E}_{\mathcal{Q}}[S|A=1] - \mathbb{E}_{\mathcal{Q}}[S|A=0], \\
\Delta_Y(\mathcal{Q}) &= \mathbb{E}_{\mathcal{Q}}[Y|A=1] - \mathbb{E}_{\mathcal{Q}}[Y|A=0].
\end{align}

Let $\phi: \mathbb{R}^2 \to \mathbb{R}$ be a functional. The estimand of interest is:
\begin{equation}
\Theta(\mu) = \mathbb{E}_{\mu}[\phi(\Delta_S(\mathcal{Q}), \Delta_Y(\mathcal{Q}))].
\end{equation}

\paragraph{Observed data.}
We observe $n$ i.i.d.\ observations $O_1, \ldots, O_n \sim \mathbb{P}_0$, where $O_i = (X_i, A_i, S_i, Y_i)$.

\paragraph{Sampling scheme.}
We draw $\mathcal{Q}_1, \ldots, \mathcal{Q}_M$ from $\mu$ via MCMC (hit-and-run sampling).

\paragraph{Estimator.}
For each sampled $\mathcal{Q}_m$, we estimate treatment effects using reweighted data:
\begin{align}
\hat{\Delta}_S^n(\mathcal{Q}_m) &= \frac{\sum_{i=1}^n w_i(\mathcal{Q}_m) A_i S_i}{\sum_{i=1}^n w_i(\mathcal{Q}_m) A_i} - \frac{\sum_{i=1}^n w_i(\mathcal{Q}_m) (1-A_i) S_i}{\sum_{i=1}^n w_i(\mathcal{Q}_m) (1-A_i)}, \\
\hat{\Delta}_Y^n(\mathcal{Q}_m) &= \frac{\sum_{i=1}^n w_i(\mathcal{Q}_m) A_i Y_i}{\sum_{i=1}^n w_i(\mathcal{Q}_m) A_i} - \frac{\sum_{i=1}^n w_i(\mathcal{Q}_m) (1-A_i) Y_i}{\sum_{i=1}^n w_i(\mathcal{Q}_m) (1-A_i)},
\end{align}
where $w_i(\mathcal{Q}) = \frac{d\mathcal{Q}}{d\mathbb{P}_0}(X_i)$ are importance weights.

The plug-in estimator is:
\begin{equation}
\hat{\Theta}_n = \frac{1}{M} \sum_{m=1}^M \phi(\hat{\Delta}_S^n(\mathcal{Q}_m), \hat{\Delta}_Y^n(\mathcal{Q}_m)).
\end{equation}

\subsection{Assumptions}

\begin{assumption}[Treatment effect estimation]
\label{assm:treatment-effects}
For each $\mathcal{Q} \in \text{supp}(\mu)$:
\begin{enumerate}[(a)]
\item Randomization or conditional exchangeability: $Y(a), S(a) \perp A \mid X$ for $a \in \{0,1\}$ under $\mathcal{Q}$.
\item Positivity: $\mathcal{Q}(A=a|X) > 0$ a.s.\ for $a \in \{0,1\}$.
\item Bounded weights: $\sup_{\mathcal{Q} \in \text{supp}(\mu)} \sup_{x \in \text{supp}(\mathbb{P}_0)} \frac{d\mathcal{Q}}{d\mathbb{P}_0}(x) \leq C < \infty$.
\item Bounded second moments: $\mathbb{E}_{\mathbb{P}_0}[S^2 + Y^2] < \infty$.
\end{enumerate}
\end{assumption}

\begin{assumption}[Functional regularity]
\label{assm:functional}
The functional $\phi: \mathbb{R}^2 \to \mathbb{R}$ is Hadamard differentiable at $(\Delta_S(\mathcal{Q}), \Delta_Y(\mathcal{Q}))$ for all $\mathcal{Q} \in \text{supp}(\mu)$, with gradient $\nabla\phi(\Delta_S, \Delta_Y) = (\partial_1 \phi, \partial_2 \phi)$ satisfying:
\begin{equation}
\sup_{\mathcal{Q} \in \text{supp}(\mu)} \|\nabla\phi(\Delta_S(\mathcal{Q}), \Delta_Y(\mathcal{Q}))\| \leq C_\phi < \infty.
\end{equation}
\end{assumption}

\begin{remark}[Aggregate functionals]
\label{rem:aggregate-gradient}
For aggregate functionals such as correlation, $\Theta(\mu) = \varphi(m(\mu))$ is a smooth function of a vector of $\mu$-moments $m(\mu)$ (five moments for correlation; see Section~\ref{sec:correlation-application}), not a pointwise map of a single pair. Throughout, $\nabla\phi(\Delta_S(\mathcal{Q}), \Delta_Y(\mathcal{Q}))$ denotes the \emph{composed} per-study gradient $\nabla\varphi(m(\mu)) \cdot \partial m/\partial(\Delta_S(\mathcal{Q}), \Delta_Y(\mathcal{Q}))$ obtained by the chain rule; the influence-function expressions \eqref{eq:aipw-s}--\eqref{eq:aipw-y} and Lemma~\ref{lem:functional-expansion} should be read with this interpretation. Non-degeneracy (Assumption~\ref{assm:nondegeneracy}) guarantees the bounded-gradient condition above for the correlation functional.
\end{remark}

\begin{assumption}[Nuisance parameter convergence]
\label{assm:nuisance}
For observational studies with estimated propensity scores $\hat{e}(X) = P(A=1|X)$ and outcome regressions $\hat{\mu}_a(X) = \mathbb{E}[Y|A=a,X]$ for $a \in \{0,1\}$, obtained via cross-fitting:
\begin{enumerate}[(a)]
\item \textbf{Strong positivity:} $e(x) \in [c, 1-c]$ uniformly for some $c \geq 0.1$.
\item \textbf{Propensity score rate:} $\|\hat{e} - e\|_{L^2(\mathbb{P}_0)} = o_P(n^{-1/4})$.
\item \textbf{Outcome regression rate:} $\|\hat{\mu}_a - \mu_a\|_{L^2(\mathbb{P}_0)} = o_P(n^{-1/4})$ for $a \in \{0,1\}$.
\item \textbf{Product convergence:} $\|\hat{e} - e\|_{L^2} \cdot \|\hat{\mu}_a - \mu_a\|_{L^2} = o_P(n^{-1/2})$.
\item \textbf{Bounded conditional moments:} $\mathbb{E}_{\mathbb{P}_0}[Y^2 | X, A=a] \leq M < \infty$ and $\mathbb{E}_{\mathbb{P}_0}[S^2 | X, A=a] \leq M < \infty$ uniformly.
\item \textbf{Weights independent of nuisances:} $w(X; \mathcal{Q}) = \frac{d\mathcal{Q}}{d\mathbb{P}_0}(X)$ does not depend on $\eta = (e, \mu_0, \mu_1)$.
\end{enumerate}
\end{assumption}

\begin{remark}[When $o_P(n^{-1/4})$ rates hold]
The $o_P(n^{-1/4})$ rates in Assumption~\ref{assm:nuisance}(b-c) hold under standard nonparametric conditions such as:
\begin{itemize}
\item Hölder smoothness $\beta > d/2$ for the nuisance functions (satisfied by GAMs, kernel methods, neural networks with appropriate architecture)
\item Sparsity or approximate sparsity assumptions (satisfied by elastic net, boosting, adaptive LASSO)
\item Structural assumptions enabling fast rates (satisfied by random forests with appropriate tuning, gradient boosting)
\end{itemize}
See Chernozhukov et al.\ (2018) for specific function class requirements. In practice, modern machine learning methods (GAMs, random forests, gradient boosting) typically achieve these rates when properly cross-validated.
\end{remark}

\begin{assumption}[Bounded second derivatives]
\label{assm:second-derivatives}
The AIPW influence function has bounded second derivatives uniformly over nuisance parameters in a neighborhood of $\eta_0 = (e, \mu_0, \mu_1)$:
\begin{equation}
\mathbb{E}_{\mathbb{P}_0}\left[\sup_{\|\eta - \eta_0\| \leq \delta} \left\|\frac{\partial^2}{\partial\eta^2}\psi_S^{\text{AIPW}}(O; \eta)\right\|^2\right] \leq B < \infty
\end{equation}
for some $\delta > 0$, and similarly for $\psi_Y^{\text{AIPW}}$.
\end{assumption}

\begin{assumption}[Uniformly bounded IF variance under reweighting]
\label{assm:bounded-if-variance}
The influence functions have uniformly bounded second moments over $\mathcal{Q} \in \text{supp}(\mu)$:
\begin{equation}
\sup_{\mathcal{Q} \in \text{supp}(\mu)} \mathbb{E}_{\mathbb{P}_0}[(w(X; \mathcal{Q}) \psi_S^{\text{AIPW}}(O; \eta_0))^2] \leq V < \infty,
\end{equation}
and similarly for $\psi_Y^{\text{AIPW}}$. This ensures the CLT applies uniformly as $M \to \infty$.
\end{assumption}

\begin{assumption}[Lipschitz gradient of functional]
\label{assm:lipschitz}
The functional $\phi: \mathbb{R}^2 \to \mathbb{R}$ has Lipschitz continuous gradient:
\begin{equation}
\|\nabla\phi(\theta_1) - \nabla\phi(\theta_2)\| \leq L\|\theta_1 - \theta_2\|
\end{equation}
for some $L < \infty$ and all $\theta_1, \theta_2$ in the range of $(\Delta_S(\mathcal{Q}), \Delta_Y(\mathcal{Q}))$ for $\mathcal{Q} \in \text{supp}(\mu)$.
\end{assumption}

\begin{assumption}[MCMC convergence]
\label{assm:mcmc}
The hit-and-run sampler satisfies geometric ergodicity with stationary distribution $\mu$. After burn-in, the empirical measure satisfies:
\begin{equation}
\left\|\hat{\mu}_M - \mu\right\|_{\text{TV}} = O_P(M^{-1/2}),
\end{equation}
where $\hat{\mu}_M = \frac{1}{M}\sum_{m=1}^M \delta_{\mathcal{Q}_m}$ and $\|\cdot\|_{\text{TV}}$ denotes total variation distance.
\end{assumption}

\begin{assumption}[Geometry regularity]
\label{assm:geometry}
The distribution $\mu$ has compact support in the weak topology: $\text{supp}(\mu)$ is compact.
\end{assumption}

\begin{assumption}[Sample size and MCMC regime]
\label{assm:sample-size}
$n \to \infty$. The primary (conditional) result treats the MCMC samples $\{\mathcal{Q}_m\}_{m=1}^M$ as fixed and holds for any $M \geq M_0$ large enough for the sampler to have reached its stationary regime. The \emph{unconditional} result, centered at the population target $\Theta(\mu)$, additionally requires $M = M(n) \to \infty$ with $n/M \to 0$ (equivalently $n = o(M)$), so that the MCMC term is asymptotically negligible relative to the $\sqrt{n}$ estimation term.
\end{assumption}

\begin{assumption}[Non-degeneracy of across-study variances]
\label{assm:nondegeneracy}
There exists $c_0 > 0$ such that $\mathrm{Var}_\mu(\Delta_S(\mathcal{Q})) \geq c_0$ and $\mathrm{Var}_\mu(\Delta_Y(\mathcal{Q})) \geq c_0$. This guarantees the correlation functional and its gradient are well-defined; the asymptotic variance scales like $c_0^{-3}$ as the boundary is approached.
\end{assumption}

\begin{remark}[MCMC sample size and the two regimes]
Decompose $\sqrt{n}(\hat{\Theta}_n - \Theta(\mu)) = \sqrt{n}(\hat{\Theta}_n - \Theta(\hat{\mu}_M)) + \sqrt{n}(\Theta(\hat{\mu}_M) - \Theta(\mu))$. The first term is $O_P(1)$ (the conditional CLT); the second is $\sqrt{n}\cdot O_P(M^{-1/2}) = O_P((n/M)^{1/2})$ by Lemma~\ref{lem:mcmc-error}. Hence:
\begin{itemize}
\item \textbf{Conditional regime (any moderate $M$):} inference targets the $M$-study correlation $\Theta(\hat{\mu}_M)$; the estimator's reported standard error is the conditional SE and is valid for $\Theta(\hat{\mu}_M)$. The residual $\Theta(\hat{\mu}_M) - \Theta(\mu) = O_P(M^{-1/2})$ is an MCMC approximation bias that does not vanish.
\item \textbf{Unconditional regime ($n = o(M)$):} the MCMC term $O_P((n/M)^{1/2}) = o_P(1)$ vanishes and $\sqrt{n}(\hat{\Theta}_n - \Theta(\mu)) \xrightarrow{d} N(0,\sigma^2(\mu))$.
\end{itemize}
Note this requires $M$ to grow \emph{faster} than $n$ (not $M = o(n)$). In the numerical studies $M \approx 2100 \gg n^{1/2}$ per effective cell, and the estimation term dominates; we report the conditional SE.
\end{remark}

\subsection{Main Result}

\begin{theorem}[Asymptotic normality, conditional on MCMC samples]
\label{thm:asymptotic-normality}
Under Assumptions~\ref{assm:treatment-effects}--\ref{assm:sample-size}, conditional on the MCMC samples $\{\mathcal{Q}_1, \ldots, \mathcal{Q}_M\}$, as $n \to \infty$:
\begin{equation}
\sqrt{n}(\hat{\Theta}_n - \Theta(\hat{\mu}_M)) \xrightarrow{d} N(0, \sigma^2(\hat{\mu}_M)),
\end{equation}
where $\Theta(\hat{\mu}_M) = \mathbb{E}_{\hat{\mu}_M}[\phi(\Delta_S(\mathcal{Q}), \Delta_Y(\mathcal{Q}))]$ and the asymptotic variance is:
\begin{equation}
\sigma^2(\hat{\mu}_M) = \mathbb{E}_{\mathbb{P}_0}[\psi_\Theta(O; \hat{\mu}_M)^2],
\end{equation}
with influence function:
\begin{equation}
\psi_\Theta(O; \hat{\mu}_M) = \mathbb{E}_{\hat{\mu}_M}\left[\nabla\phi(\Delta_S(\mathcal{Q}), \Delta_Y(\mathcal{Q})) \cdot \begin{pmatrix} \psi_S^{\text{AIPW}}(O; \eta_0, \mathcal{Q}) \\ \psi_Y^{\text{AIPW}}(O; \eta_0, \mathcal{Q}) \end{pmatrix}\right],
\end{equation}
where $\psi_S^{\text{AIPW}}(O; \eta, \mathcal{Q})$ and $\psi_Y^{\text{AIPW}}(O; \eta, \mathcal{Q})$ are the AIPW influence functions defined in \eqref{eq:aipw-s}--\eqref{eq:aipw-y}. For randomized trials, these reduce to weighted difference-in-means; for observational studies with cross-fitting, they provide doubly robust estimation.

Furthermore:
\begin{enumerate}[(a)]
\item If $M$ is fixed (or grows slower than $n$), $\Theta(\hat{\mu}_M) - \Theta(\mu) = O_P(M^{-1/2})$ (MCMC approximation bias), and inference is for $\Theta(\hat{\mu}_M)$.
\item If $M = M(n) \to \infty$ with $n/M \to 0$, the MCMC term $\sqrt{n}(\Theta(\hat{\mu}_M) - \Theta(\mu)) = O_P((n/M)^{1/2}) = o_P(1)$, and unconditionally:
\begin{equation}
\sqrt{n}(\hat{\Theta}_n - \Theta(\mu)) \xrightarrow{d} N(0, \sigma^2(\mu)).
\end{equation}
\end{enumerate}
\end{theorem}

\begin{remark}[Practical interpretation]
Conditioning on MCMC samples separates two sources of randomness:
\begin{itemize}
\item \textbf{Sampling variation} (from finite $n$): Captured by the asymptotic distribution, rate $n^{-1/2}$.
\item \textbf{MCMC approximation} (from finite $M$): Bias in target, magnitude $M^{-1/2}$.
\end{itemize}
When $M$ is large relative to $n$, the MCMC bias $O_P(M^{-1/2})$ is negligible relative to the $O_P(n^{-1/2})$ sampling variation; formally it vanishes in the regime $n = o(M)$.
\end{remark}

\subsection{Proof of Theorem~\ref{thm:asymptotic-normality}}

\subsubsection{Proof roadmap}

The proof proceeds in four main steps:
\begin{enumerate}
\item \textbf{MCMC approximation:} Show (via a Markov chain CLT) that the empirical measure over MCMC samples approximates $\mu$ at rate $M^{-1/2}$. Relative to the $\sqrt{n}$ estimation term this contributes $O_P((n/M)^{1/2})$, negligible only in the unconditional regime $n = o(M)$ (Assumption~\ref{assm:sample-size}); otherwise inference is conditional on $\hat{\mu}_M$.
\item \textbf{Treatment effect expansion (AIPW):} Establish Neyman orthogonality under reweighting (Lemma~\ref{lem:orthogonality}), verify conditions for cross-fitting with bounded weights (Lemma~\ref{lem:cross-fitting}), and obtain uniform influence function expansion with $o_P(n^{-1/2})$ remainder. Key innovation: bounded second derivatives + bounded weights give deterministic uniformity over MCMC samples.
\item \textbf{Functional delta method:} Apply Hadamard differentiability to obtain the expansion for $\phi(\hat{\Delta}_S^n(\mathcal{Q}), \hat{\Delta}_Y^n(\mathcal{Q}))$.
\item \textbf{Integration and CLT:} Average over $\mu$ (or its MCMC approximation) and apply the central limit theorem using uniform moment bounds.
\end{enumerate}

\subsubsection{Step 1: MCMC approximation}

Let $\hat{\mu}_M = \frac{1}{M}\sum_{m=1}^M \delta_{\mathcal{Q}_m}$ denote the empirical measure over MCMC samples.

\begin{lemma}[MCMC approximation error]
\label{lem:mcmc-error}
Under Assumptions~\ref{assm:mcmc} (geometric ergodicity) and \ref{assm:geometry} (compact support), for any bounded measurable function $g: \text{supp}(\mu) \to \mathbb{R}$ with $\|g\|_\infty \leq 1$:
\begin{equation}
\left|\mathbb{E}_{\hat{\mu}_M}[g(\mathcal{Q})] - \mathbb{E}_{\mu}[g(\mathcal{Q})]\right| = O_P(M^{-1/2}).
\end{equation}
\end{lemma}

\begin{proof}[Proof of Lemma~\ref{lem:mcmc-error}]
The hit-and-run samples $\mathcal{Q}_1, \ldots, \mathcal{Q}_M$ form a \emph{dependent} Markov chain, not independent draws; we therefore invoke the Markov chain CLT rather than the i.i.d.\ CLT. By Assumption~\ref{assm:mcmc}, the chain is geometrically ergodic with stationary distribution $\mu$, and $g$ is bounded (hence in $L^2(\mu)$ with all polynomial moments). Under geometric ergodicity and a bounded (thus $L^{2+\delta}(\mu)$) functional, the Markov chain CLT holds \cite{jones2004markov, roberts2004general}:
\begin{equation}
\sqrt{M}\left(\frac{1}{M}\sum_{m=1}^M g(\mathcal{Q}_m) - \mathbb{E}_{\mu}[g]\right) \xrightarrow{d} N\!\left(0, \sigma_g^2\right),
\end{equation}
where $\sigma_g^2 = \mathrm{Var}_\mu(g(\mathcal{Q}_1)) + 2\sum_{k=1}^\infty \mathrm{Cov}_\mu(g(\mathcal{Q}_1), g(\mathcal{Q}_{1+k}))$ is the \emph{integrated-autocorrelation} asymptotic variance (finite by geometric ergodicity). Hence $\mathbb{E}_{\hat{\mu}_M}[g] - \mathbb{E}_{\mu}[g] = O_P(M^{-1/2})$. The rate is unchanged from the independent case; only the variance constant differs (it inflates by the integrated autocorrelation time). Applying this to $g = \phi(\Delta_S(\cdot), \Delta_Y(\cdot))$ (bounded under Assumptions~\ref{assm:treatment-effects} and \ref{assm:functional}) gives $\Theta(\hat{\mu}_M) - \Theta(\mu) = O_P(M^{-1/2})$.
\end{proof}

\begin{corollary}
\label{cor:mcmc-negligible}
Under Assumption~\ref{assm:sample-size}:
\begin{enumerate}[(a)]
\item If $M$ is fixed, $|\mathbb{E}_{\hat{\mu}_M}[g] - \mathbb{E}_{\mu}[g]| = O_P(M^{-1/2})$ is constant in $n$.
\item If $M \to \infty$, $|\mathbb{E}_{\hat{\mu}_M}[g] - \mathbb{E}_{\mu}[g]| = o_P(1)$ vanishes as $n \to \infty$.
\end{enumerate}
\end{corollary}

\begin{remark}
In the fixed-$M$ case, the theorem establishes asymptotic normality for $\Theta(\hat{\mu}_M)$ (the MCMC-approximated target) rather than $\Theta(\mu)$ itself. The difference $|\Theta(\hat{\mu}_M) - \Theta(\mu)| = O_P(M^{-1/2})$ is a source of bias that does not vanish unless $M \to \infty$. It is negligible relative to $\sqrt{n}$ sampling variation precisely in the regime $n = o(M)$.
\end{remark}

\subsubsection{Step 2: Treatment effect expansion}

\paragraph{AIPW influence functions.}
For observational studies with estimated nuisances, we use augmented inverse probability weighting (AIPW). The influence functions are:
\begin{align}
\psi_S^{\text{AIPW}}(O; \eta, \mathcal{Q}) &= w(X; \mathcal{Q}) \Bigg[\frac{A(S - \mu_1^S(X))}{e(X)} - \frac{(1-A)(S - \mu_0^S(X))}{1-e(X)} \nonumber \\
&\quad + \mu_1^S(X) - \mu_0^S(X)\Bigg] - \Delta_S(\mathcal{Q}), \label{eq:aipw-s} \\
\psi_Y^{\text{AIPW}}(O; \eta, \mathcal{Q}) &= w(X; \mathcal{Q}) \Bigg[\frac{A(Y - \mu_1^Y(X))}{e(X)} - \frac{(1-A)(Y - \mu_0^Y(X))}{1-e(X)} \nonumber \\
&\quad + \mu_1^Y(X) - \mu_0^Y(X)\Bigg] - \Delta_Y(\mathcal{Q}), \label{eq:aipw-y}
\end{align}
where:
\begin{itemize}
\item $\eta = (e, \mu_0^S, \mu_1^S, \mu_0^Y, \mu_1^Y)$ are nuisance parameters,
\item $e(X) = P(A=1|X)$ is the propensity score,
\item $\mu_a^S(X) = \mathbb{E}[S|A=a,X]$ and $\mu_a^Y(X) = \mathbb{E}[Y|A=a,X]$ are outcome regressions,
\item $w(X; \mathcal{Q}) = \frac{d\mathcal{Q}}{d\mathbb{P}_0}(X)$ are importance weights (independent of $\eta$ by Assumption~\ref{assm:nuisance}(f)),
\item $\Delta_S(\mathcal{Q}) = \mathbb{E}_{\mathcal{Q}}[\mu_1^S(X) - \mu_0^S(X)]$ and $\Delta_Y(\mathcal{Q}) = \mathbb{E}_{\mathcal{Q}}[\mu_1^Y(X) - \mu_0^Y(X)]$.
\end{itemize}

\begin{remark}[Three-term decomposition]
The AIPW influence function has three terms:
\begin{enumerate}
\item \textbf{IPW term:} $\frac{A(S - \mu_1^S(X))}{e(X)} - \frac{(1-A)(S - \mu_0^S(X))}{1-e(X)}$ (inverse probability weighting of residuals)
\item \textbf{Regression term:} $\mu_1^S(X) - \mu_0^S(X)$ (predicted treatment effect)
\item \textbf{Centering:} the \emph{constant} $-\Delta_S(\mathcal{Q})$, subtracted \emph{outside} the importance weight, ensuring $\mathbb{E}_{\mathbb{P}_0}[\psi_S^{\text{AIPW}}] = 0$.
\end{enumerate}
This structure provides \textbf{double robustness}: the estimator is consistent if either the propensity score $e(X)$ or the outcome regressions $\mu_a(X)$ are correctly specified (not both required).

Two features are specific to this estimand and distinguish it from the textbook difference-in-means AIPW score. First, the entire bracketed AIPW score, \emph{including} the regression term $\mu_1^S(X) - \mu_0^S(X)$, is multiplied by the importance weight $w(X;\mathcal{Q})$. This is correct because the target $\Delta_S(\mathcal{Q}) = \mathbb{E}_{\mathcal{Q}}[\mu_1^S(X) - \mu_0^S(X)]$ is an expectation under the \emph{fixed, external} composition $\mathcal{Q}$, not under $\mathbb{P}_0$; there is consequently no $\mathbb{P}_0$-covariate-law augmentation term. Second, the centering is the constant $-\Delta_S(\mathcal{Q})$, not $-w(X;\mathcal{Q})\Delta_S(\mathcal{Q})$: only constant centering yields a mean-zero per-observation influence function with the correct variance (see the companion derivation, \texttt{derivation\_influence\_functions.md}).
\end{remark}

\begin{remark}[Special case: RCTs]
For randomized controlled trials where $A \perp (S, Y, X)$ under $\mathbb{P}_0$, the propensity score $e(X) = P(A=1)$ is known and constant. Setting $\mu_a^S(X) = \mathbb{E}[S|A=a]$ and $\mu_a^Y(X) = \mathbb{E}[Y|A=a]$ (constants), the AIPW IF reduces to the simple weighted difference-in-means IF used in the RCT-only proof. Thus, the AIPW framework generalizes the RCT case.
\end{remark}

\begin{lemma}[Neyman orthogonality under reweighting]
\label{lem:orthogonality}
Under Assumptions~\ref{assm:treatment-effects} and \ref{assm:nuisance}, for the weighted AIPW functional $\theta_w(\eta; \mathcal{Q}) = \mathbb{E}_{\mathbb{P}_0}[w(X; \mathcal{Q}) \psi_S^{\text{AIPW}}(O; \eta, \mathcal{Q})]$:
\begin{equation}
\frac{\partial}{\partial\eta}\theta_w(\eta; \mathcal{Q})\bigg|_{\eta = \eta_0}[r] = 0
\end{equation}
for all $r$ in the tangent space and all $\mathcal{Q} \in \text{supp}(\mu)$.
\end{lemma}

\begin{proof}[Proof of Lemma~\ref{lem:orthogonality}]
We prove the result for $\psi_S^{\text{AIPW}}$; the proof for $\psi_Y^{\text{AIPW}}$ is identical.

\paragraph{Step 1: Decompose AIPW score.}
Write the AIPW score as:
\begin{equation}
\psi_S^{\text{AIPW}}(O; \eta, \mathcal{Q}) = \psi_S^{\text{IPW}}(O; \eta) + \psi_S^{\text{reg}}(O; \eta) - \Delta_S(\mathcal{Q}),
\end{equation}
where:
\begin{align}
\psi_S^{\text{IPW}}(O; \eta) &= \frac{A(S - \mu_1^S(X))}{e(X)} - \frac{(1-A)(S - \mu_0^S(X))}{1-e(X)}, \\
\psi_S^{\text{reg}}(O; \eta) &= \mu_1^S(X) - \mu_0^S(X).
\end{align}

\paragraph{Step 2: Standard AIPW orthogonality (conditional).}
For standard AIPW (without reweighting), Neyman orthogonality is well-established (Chernozhukov et al., 2018): the conditional expectation of the pathwise derivative vanishes:
\begin{equation}
\mathbb{E}\left[\frac{\partial}{\partial\eta}\psi_S^{\text{AIPW}}(O; \eta_0)[r] \,\bigg|\, X\right] = 0
\end{equation}
for all $r$ in the tangent space. This is the \textbf{conditional mean-zero property}.

\paragraph{Step 3: Weights don't depend on nuisances.}
By Assumption~\ref{assm:nuisance}(f), $w(X; \mathcal{Q}) = \frac{d\mathcal{Q}}{d\mathbb{P}_0}(X)$ is a function of $X$ only and does not depend on $\eta = (e, \mu_0, \mu_1)$.

\paragraph{Step 4: Orthogonality survives weighting.}
For the weighted functional:
\begin{align}
\frac{\partial}{\partial\eta}\theta_w(\eta_0; \mathcal{Q})[r] &= \frac{\partial}{\partial\eta}\mathbb{E}_{\mathbb{P}_0}\left[w(X; \mathcal{Q}) \psi_S^{\text{AIPW}}(O; \eta_0, \mathcal{Q})\right][r] \\
&= \mathbb{E}_{\mathbb{P}_0}\left[w(X; \mathcal{Q}) \cdot \frac{\partial}{\partial\eta}\psi_S^{\text{AIPW}}(O; \eta_0, \mathcal{Q})[r]\right] \\
&= \mathbb{E}_{\mathbb{P}_0}\left[w(X; \mathcal{Q}) \cdot \mathbb{E}\left[\frac{\partial}{\partial\eta}\psi_S^{\text{AIPW}}(O; \eta_0)[r] \,\bigg|\, X\right]\right] \\
&= \mathbb{E}_{\mathbb{P}_0}[w(X; \mathcal{Q}) \cdot 0] = 0.
\end{align}
The third equality uses the law of iterated expectations, and the fourth uses the conditional mean-zero property from Step 2.

\paragraph{Conclusion.}
Neyman orthogonality holds for the weighted functional. This is the key property that eliminates first-order nuisance bias.
\end{proof}

\begin{lemma}[Cross-fitting with bounded weights]
\label{lem:cross-fitting}
Under Assumptions~\ref{assm:treatment-effects}, \ref{assm:nuisance}, and \ref{assm:second-derivatives}, the cross-fitted AIPW estimator with bounded weights satisfies:
\begin{equation}
\hat{\Delta}_S(\mathcal{Q}) - \Delta_S(\mathcal{Q}) = \frac{1}{n}\sum_{i=1}^n w_i(\mathcal{Q}) \psi_S^{\text{AIPW}}(O_i; \eta_0, \mathcal{Q}) + o_P(n^{-1/2}),
\end{equation}
uniformly over $\mathcal{Q} \in \text{supp}(\mu)$. Similarly for $\hat{\Delta}_Y(\mathcal{Q})$.
\end{lemma}

\begin{proof}[Proof of Lemma~\ref{lem:cross-fitting}]
We verify the conditions of Theorem 3.1 in Chernozhukov et al.\ (2018) for the weighted AIPW estimator.

\paragraph{Condition 1: Neyman orthogonality.}
Verified in Lemma~\ref{lem:orthogonality}. The weighted functional $\theta_w(\eta; \mathcal{Q})$ satisfies $\frac{\partial}{\partial\eta}\theta_w(\eta_0; \mathcal{Q}) = 0$.

\paragraph{Condition 2: $L^2$ convergence rates.}
By Assumption~\ref{assm:nuisance}(b-c):
\begin{align}
\|\hat{e} - e\|_{L^2(\mathbb{P}_0)} &= o_P(n^{-1/4}), \\
\|\hat{\mu}_a - \mu_a\|_{L^2(\mathbb{P}_0)} &= o_P(n^{-1/4}) \quad \text{for } a \in \{0,1\}.
\end{align}

\paragraph{Condition 3: Bounded moments.}
By Assumption~\ref{assm:nuisance}(e), $\mathbb{E}[Y^2|X,A], \mathbb{E}[S^2|X,A] \leq M < \infty$, ensuring the influence function has finite second moments.

\paragraph{Condition 4: Cross-fitting.}
Sample splitting into $K$ folds ensures nuisance estimators $\hat{\eta}^{(-k)}$ are independent of observations in fold $k$, eliminating overfitting bias.

\paragraph{Condition 5: Bounded weights preserve all properties.}
By Assumption~\ref{assm:treatment-effects}(c), $w_i(\mathcal{Q}) \leq C < \infty$ uniformly. Bounded weights preserve:
\begin{itemize}
\item Neyman orthogonality (Lemma~\ref{lem:orthogonality})
\item $L^2$ convergence of weighted IF: $\|w \cdot \psi^{\text{AIPW}}(\cdot; \hat{\eta}) - w \cdot \psi^{\text{AIPW}}(\cdot; \eta_0)\|_{L^2} \leq C \|\psi^{\text{AIPW}}(\cdot; \hat{\eta}) - \psi^{\text{AIPW}}(\cdot; \eta_0)\|_{L^2}$
\item Bounded second moments (Assumption~\ref{assm:bounded-if-variance})
\end{itemize}

\paragraph{Second-order remainder bound.}
By Assumption~\ref{assm:second-derivatives}, second derivatives of $\psi_S^{\text{AIPW}}$ are bounded. Taylor expansion gives:
\begin{align}
&\hat{\Delta}_S(\mathcal{Q}) - \Delta_S(\mathcal{Q}) \\
&= \frac{1}{n}\sum_{i=1}^n w_i(\mathcal{Q}) \psi_S^{\text{AIPW}}(O_i; \eta_0, \mathcal{Q}) + \frac{1}{n}\sum_{i=1}^n w_i(\mathcal{Q}) R_i(\hat{\eta}, \eta_0),
\end{align}
where the remainder satisfies:
\begin{equation}
R_i(\hat{\eta}, \eta_0) = \int_0^1 (1-t) \frac{\partial^2}{\partial\eta^2}\psi_S^{\text{AIPW}}(O_i; \tilde{\eta}_t, \mathcal{Q})[\hat{\eta} - \eta_0]^2 dt,
\end{equation}
with $\tilde{\eta}_t = \eta_0 + t(\hat{\eta} - \eta_0)$.

By Assumption~\ref{assm:second-derivatives} and bounded weights:
\begin{equation}
\left|\frac{1}{n}\sum_{i=1}^n w_i(\mathcal{Q}) R_i\right| \leq C \cdot B \cdot \|\hat{\eta} - \eta_0\|_{L^2}^2 = C \cdot B \cdot o_P(n^{-1/2}),
\end{equation}
where the last equality uses Assumption~\ref{assm:nuisance}(d).

\paragraph{Uniformity over $\mathcal{Q}$.}
The remainder bound $|R_n(\mathcal{Q})| \leq CB\|\hat{\eta} - \eta_0\|_{L^2}^2$ is \textbf{independent of $\mathcal{Q}$} (depends only on $C, B, \|\hat{\eta} - \eta_0\|_{L^2}$). Therefore:
\begin{equation}
\sup_{\mathcal{Q} \in \text{supp}(\mu)} |R_n(\mathcal{Q})| = o_P(n^{-1/2}).
\end{equation}
This is \textbf{deterministic uniformity}—no maximal inequalities or Donsker theory required.

\paragraph{Conclusion.}
By Chernozhukov et al.\ (2018, Theorem 3.1), the expansion holds with $o_P(n^{-1/2})$ remainder, uniformly over $\mathcal{Q}$.
\end{proof}


\subsubsection{Step 3: Functional delta method}

\begin{lemma}[Functional expansion]
\label{lem:functional-expansion}
Under Assumptions~\ref{assm:treatment-effects}, \ref{assm:nuisance}, \ref{assm:functional}, and \ref{assm:lipschitz}, uniformly over $\mathcal{Q} \in \text{supp}(\mu)$:
\begin{align}
&\sqrt{n}(\phi(\hat{\Delta}_S^n(\mathcal{Q}), \hat{\Delta}_Y^n(\mathcal{Q})) - \phi(\Delta_S(\mathcal{Q}), \Delta_Y(\mathcal{Q}))) \nonumber \\
&\quad = \nabla\phi(\Delta_S(\mathcal{Q}), \Delta_Y(\mathcal{Q})) \cdot \begin{pmatrix} \frac{1}{\sqrt{n}}\sum_{i=1}^n \psi_S^{\text{AIPW}}(O_i; \eta_0, \mathcal{Q}) \\ \frac{1}{\sqrt{n}}\sum_{i=1}^n \psi_Y^{\text{AIPW}}(O_i; \eta_0, \mathcal{Q}) \end{pmatrix} + o_P(1).
\end{align}
\end{lemma}

\begin{proof}[Proof of Lemma~\ref{lem:functional-expansion}]
By Hadamard differentiability of $\phi$ (Assumption~\ref{assm:functional}):
\begin{align}
&\phi(\hat{\Delta}_S^n(\mathcal{Q}), \hat{\Delta}_Y^n(\mathcal{Q})) - \phi(\Delta_S(\mathcal{Q}), \Delta_Y(\mathcal{Q})) \\
&\quad = \nabla\phi(\Delta_S(\mathcal{Q}), \Delta_Y(\mathcal{Q})) \cdot \begin{pmatrix} \hat{\Delta}_S^n(\mathcal{Q}) - \Delta_S(\mathcal{Q}) \\ \hat{\Delta}_Y^n(\mathcal{Q}) - \Delta_Y(\mathcal{Q}) \end{pmatrix} + o(\|(\hat{\Delta}_S^n, \hat{\Delta}_Y^n) - (\Delta_S, \Delta_Y)\|).
\end{align}

By Lemma~\ref{lem:cross-fitting}, $\|(\hat{\Delta}_S^n, \hat{\Delta}_Y^n) - (\Delta_S, \Delta_Y)\| = O_P(n^{-1/2})$, hence the remainder is $o_P(n^{-1/2})$. Substituting the expansion from Lemma~\ref{lem:cross-fitting} gives the result.
\end{proof}

\subsubsection{Step 4: Integration and CLT}

\begin{proof}[Proof of Theorem~\ref{thm:asymptotic-normality}]
We combine the preceding lemmas in four steps.

\paragraph{Step 4.1: Condition on MCMC samples.}
Throughout this proof, we condition on the MCMC samples $\{\mathcal{Q}_1, \ldots, \mathcal{Q}_M\}$, treating $\hat{\mu}_M$ as fixed. The estimator becomes:
\begin{equation}
\hat{\Theta}_n = \mathbb{E}_{\hat{\mu}_M}[\phi(\hat{\Delta}_S^n(\mathcal{Q}), \hat{\Delta}_Y^n(\mathcal{Q}))] = \frac{1}{M}\sum_{m=1}^M \phi(\hat{\Delta}_S^n(\mathcal{Q}_m), \hat{\Delta}_Y^n(\mathcal{Q}_m)),
\end{equation}
and the target is:
\begin{equation}
\Theta(\hat{\mu}_M) = \mathbb{E}_{\hat{\mu}_M}[\phi(\Delta_S(\mathcal{Q}), \Delta_Y(\mathcal{Q}))].
\end{equation}

\paragraph{Step 4.2: Decompose estimator (conditional).}
\begin{equation}
\sqrt{n}(\hat{\Theta}_n - \Theta(\hat{\mu}_M)) = \sqrt{n}\left(\mathbb{E}_{\hat{\mu}_M}[\phi(\hat{\Delta}_S^n, \hat{\Delta}_Y^n)] - \mathbb{E}_{\hat{\mu}_M}[\phi(\Delta_S, \Delta_Y)]\right).
\end{equation}

\paragraph{Step 4.3: Expand using functional delta method.}
By Lemma~\ref{lem:functional-expansion}:
\begin{align}
\sqrt{n}(\hat{\Theta}_n - \Theta(\hat{\mu}_M)) &= \sqrt{n} \mathbb{E}_{\hat{\mu}_M}\left[\nabla\phi(\Delta_S, \Delta_Y) \cdot \begin{pmatrix} \hat{\Delta}_S^n - \Delta_S \\ \hat{\Delta}_Y^n - \Delta_Y \end{pmatrix}\right] + o_P(1) \\
&= \mathbb{E}_{\hat{\mu}_M}\left[\nabla\phi(\Delta_S, \Delta_Y) \cdot \begin{pmatrix} \frac{1}{\sqrt{n}}\sum_{i=1}^n \psi_S^{\text{AIPW}}(O_i; \eta_0, \mathcal{Q}) \\ \frac{1}{\sqrt{n}}\sum_{i=1}^n \psi_Y^{\text{AIPW}}(O_i; \eta_0, \mathcal{Q}) \end{pmatrix}\right] + o_P(1).
\end{align}

\paragraph{Step 4.4: Interchange expectation and sum.}
Since $\hat{\mu}_M$ is treated as fixed (conditional on MCMC samples), and by Fubini's theorem:
\begin{align}
\sqrt{n}(\hat{\Theta}_n - \Theta(\hat{\mu}_M)) &= \frac{1}{\sqrt{n}}\sum_{i=1}^n \mathbb{E}_{\hat{\mu}_M}\left[\nabla\phi(\Delta_S, \Delta_Y) \cdot \begin{pmatrix} \psi_S^{\text{AIPW}}(O_i; \eta_0, \mathcal{Q}) \\ \psi_Y^{\text{AIPW}}(O_i; \eta_0, \mathcal{Q}) \end{pmatrix}\right] + o_P(1) \\
&= \frac{1}{\sqrt{n}}\sum_{i=1}^n \psi_\Theta(O_i; \hat{\mu}_M) + o_P(1),
\end{align}
where the composite influence function is:
\begin{equation}
\psi_\Theta(O; \hat{\mu}_M) = \mathbb{E}_{\hat{\mu}_M}\left[\nabla\phi(\Delta_S(\mathcal{Q}), \Delta_Y(\mathcal{Q})) \cdot \begin{pmatrix} \psi_S^{\text{AIPW}}(O; \eta_0, \mathcal{Q}) \\ \psi_Y^{\text{AIPW}}(O; \eta_0, \mathcal{Q}) \end{pmatrix}\right].
\end{equation}

\paragraph{Step 4.5: Apply CLT (conditional).}
Conditional on $\{\mathcal{Q}_1, \ldots, \mathcal{Q}_M\}$, the observations $O_1, \ldots, O_n \overset{\text{i.i.d.}}{\sim} \mathbb{P}_0$ are independent of the MCMC samples. By Assumption~\ref{assm:bounded-if-variance}, $\mathbb{E}_{\mathbb{P}_0}[\psi_\Theta(O; \hat{\mu}_M)^2] < \infty$. By the CLT:
\begin{equation}
\frac{1}{\sqrt{n}}\sum_{i=1}^n \psi_\Theta(O_i; \hat{\mu}_M) \xrightarrow{d} N(0, \mathbb{E}_{\mathbb{P}_0}[\psi_\Theta(O; \hat{\mu}_M)^2]) = N(0, \sigma^2(\hat{\mu}_M)).
\end{equation}

Note that $\mathbb{E}_{\mathbb{P}_0}[\psi_\Theta(O; \hat{\mu}_M)] = 0$ by construction (using Neyman orthogonality from Lemma~\ref{lem:orthogonality}), conditional on the MCMC samples.

\paragraph{Step 4.6: Unconditional result (when $M \to \infty$).}
When $M \to \infty$, by Assumption~\ref{assm:mcmc}, $\hat{\mu}_M \to \mu$ weakly. By Assumption~\ref{assm:bounded-if-variance} (uniform moment bound), $\sup_{\mathcal{Q}} \mathbb{E}[\psi^{\text{AIPW}}(O; \mathcal{Q})^2] < \infty$, hence $\psi_\Theta(O; \hat{\mu}_M) \to \psi_\Theta(O; \mu)$ in $L^2(\mathbb{P}_0)$ by dominated convergence (using bounded $\nabla\phi$ from Assumption~\ref{assm:functional}). Therefore:
\begin{equation}
\sigma^2(\hat{\mu}_M) \to \sigma^2(\mu),
\end{equation}
and unconditionally:
\begin{equation}
\sqrt{n}(\hat{\Theta}_n - \Theta(\mu)) \xrightarrow{d} N(0, \sigma^2(\mu)).
\end{equation}

\paragraph{Conclusion.}
We have shown:
\begin{enumerate}[(a)]
\item Conditional on MCMC samples: $\sqrt{n}(\hat{\Theta}_n - \Theta(\hat{\mu}_M)) \xrightarrow{d} N(0, \sigma^2(\hat{\mu}_M))$.
\item When $M \to \infty$: $\Theta(\hat{\mu}_M) \to \Theta(\mu)$ and $\sigma^2(\hat{\mu}_M) \to \sigma^2(\mu)$.
\end{enumerate}
\end{proof}

\subsection{Variance Estimation}

The asymptotic variance $\sigma^2(\hat{\mu}_M)$ can be consistently estimated (conditional on MCMC samples) by:
\begin{equation}
\hat{\sigma}^2_n = \frac{1}{n}\sum_{i=1}^n \hat{\psi}_\Theta(O_i; \hat{\mu}_M)^2,
\end{equation}
where:
\begin{equation}
\hat{\psi}_\Theta(O_i; \hat{\mu}_M) = \frac{1}{M}\sum_{m=1}^M \nabla\phi(\hat{\Delta}_S^n(\mathcal{Q}_m), \hat{\Delta}_Y^n(\mathcal{Q}_m)) \cdot \begin{pmatrix} \psi_S^{\text{AIPW}}(O_i; \hat{\eta}^{(-k_i)}, \mathcal{Q}_m) \\ \psi_Y^{\text{AIPW}}(O_i; \hat{\eta}^{(-k_i)}, \mathcal{Q}_m) \end{pmatrix}.
\end{equation}
Here, $\hat{\eta}^{(-k_i)}$ denotes the nuisance parameters estimated on the cross-fitting fold that excludes observation $i$. This ensures the influence function estimates are independent of the data used to estimate nuisances, maintaining the double robustness property.

\begin{proposition}[Consistent variance estimation]
\label{prop:variance-consistency}
Under the conditions of Theorem~\ref{thm:asymptotic-normality}, conditional on $\{\mathcal{Q}_1, \ldots, \mathcal{Q}_M\}$:
\begin{equation}
\hat{\sigma}^2_n \xrightarrow{P} \sigma^2(\hat{\mu}_M) \quad \text{as } n \to \infty.
\end{equation}
Furthermore, if $M \to \infty$, then unconditionally $\hat{\sigma}^2_n \xrightarrow{P} \sigma^2(\mu)$.
\end{proposition}

\begin{proof}[Proof of Proposition~\ref{prop:variance-consistency}]
Consistency follows in three steps, plus the unconditional extension.

\paragraph{Step 1: $L^2$ convergence of AIPW influence functions.}
By Assumption~\ref{assm:nuisance} (nuisance convergence rates) and the Lipschitz structure of the AIPW score:
\begin{equation}
\|\psi_S^{\text{AIPW}}(O; \hat{\eta}, \mathcal{Q}) - \psi_S^{\text{AIPW}}(O; \eta_0, \mathcal{Q})\|_{L^2(\mathbb{P}_0)} \xrightarrow{P} 0.
\end{equation}
This follows from:
\begin{itemize}
\item $\|\hat{e} - e\|_{L^2} = o_P(1)$ and $\|\hat{\mu}_a - \mu_a\|_{L^2} = o_P(1)$ (Assumption~\ref{assm:nuisance}(b-c))
\item AIPW score is Lipschitz in $(e, \mu_0, \mu_1)$ on the domain bounded away from $0$ and $1$ (by positivity, Assumption~\ref{assm:nuisance}(a))
\item Bounded weights: $w(X; \mathcal{Q}) \leq C$ (Assumption~\ref{assm:treatment-effects}(c))
\end{itemize}

\paragraph{Step 2: Gradient convergence.}
By Lemma~\ref{lem:cross-fitting}, $\hat{\Delta}_S(\mathcal{Q}_m) \xrightarrow{P} \Delta_S(\mathcal{Q}_m)$ and $\hat{\Delta}_Y(\mathcal{Q}_m) \xrightarrow{P} \Delta_Y(\mathcal{Q}_m)$ for each $m$. By Assumption~\ref{assm:lipschitz} (Lipschitz gradient):
\begin{equation}
\nabla\phi(\hat{\Delta}_S(\mathcal{Q}_m), \hat{\Delta}_Y(\mathcal{Q}_m)) \xrightarrow{P} \nabla\phi(\Delta_S(\mathcal{Q}_m), \Delta_Y(\mathcal{Q}_m)).
\end{equation}

\paragraph{Step 3: Combine via continuous mapping.}
Define $g: L^2(\mathbb{P}_0) \times \mathbb{R}^2 \to \mathbb{R}$ by:
\begin{equation}
g(\psi, \nabla) = \mathbb{E}_{\mathbb{P}_0}[(\nabla \cdot \psi)^2].
\end{equation}
By Steps 1-2, $(\hat{\psi}^{\text{AIPW}}, \nabla\phi(\hat{\Delta})) \xrightarrow{P} (\psi^{\text{AIPW}}, \nabla\phi(\Delta))$ in the product topology. Since $g$ is continuous (by dominated convergence using bounded moments), the continuous mapping theorem gives:
\begin{equation}
\hat{\sigma}^2_n = \frac{1}{n}\sum_{i=1}^n \hat{\psi}_\Theta(O_i)^2 \xrightarrow{P} \mathbb{E}_{\mathbb{P}_0}[\psi_\Theta(O)^2] = \sigma^2(\hat{\mu}_M).
\end{equation}

\paragraph{Step 4: Unconditional convergence when $M \to \infty$.}
When $M \to \infty$, by Assumption~\ref{assm:mcmc}, $\hat{\mu}_M \to \mu$ weakly. By Assumption~\ref{assm:bounded-if-variance} (uniform moment bound), dominated convergence applies, giving $\sigma^2(\hat{\mu}_M) \to \sigma^2(\mu)$. Combining with Step 3 yields unconditional consistency.
\end{proof}

\begin{remark}[Practical variance estimation]
In practice, one estimates $\sigma^2(\hat{\mu}_M)$, which accounts for sampling variation but not MCMC approximation error. If $M$ is large (e.g., $M \geq 500$), the MCMC error is typically negligible relative to sampling uncertainty. Formal accounting for MCMC error would require bootstrap or other resampling methods applied to the MCMC procedure itself, but this is rarely done in practice.
\end{remark}

\subsection{Application to Correlation Functional}\label{sec:correlation-application}

For the correlation functional $\phi(\Delta_S, \Delta_Y) = \text{cor}_\mu(\Delta_S, \Delta_Y)$, this is a composite functional:
\begin{equation}
\phi(\Delta_S, \Delta_Y) = \frac{\mathbb{E}_\mu[\Delta_S \Delta_Y] - \mathbb{E}_\mu[\Delta_S]\mathbb{E}_\mu[\Delta_Y]}{\sqrt{\text{Var}_\mu(\Delta_S) \text{Var}_\mu(\Delta_Y)}}.
\end{equation}

Correlation is thus a smooth function $\varphi$ of the \emph{five base moments} $m(\mu) = (\mathbb{E}_\mu[\Delta_S], \mathbb{E}_\mu[\Delta_Y], \mathbb{E}_\mu[\Delta_S^2], \mathbb{E}_\mu[\Delta_Y^2], \mathbb{E}_\mu[\Delta_S \Delta_Y])$, i.e. $\Theta(\mu) = \varphi(m(\mu))$. This is the object to which the functional delta method applies: $\phi$ is not a fixed map $\mathbb{R}^2 \to \mathbb{R}$ of a single pair $(\Delta_S, \Delta_Y)$, but a functional of the across-study \emph{distribution} $\mu$ of the pair, expressed through these five $\mu$-moments. Each base moment is a smooth (polynomial) map of $(\Delta_S(\mathcal{Q}), \Delta_Y(\mathcal{Q}))$ averaged over $\mathcal{Q} \sim \mu$, so the per-study gradient $\nabla\phi(\Delta_S(\mathcal{Q}_m), \Delta_Y(\mathcal{Q}_m))$ used in \eqref{eq:aipw-s}--\eqref{eq:aipw-y} and Lemma~\ref{lem:functional-expansion} is the chain-rule composition of $\nabla\varphi(m(\mu))$ with the moment derivatives.

\begin{remark}
The map $\varphi$ is continuously differentiable, hence Hadamard differentiable, on the region $\{\mathrm{Var}_\mu(\Delta_S) \geq c_0,\ \mathrm{Var}_\mu(\Delta_Y) \geq c_0\}$ guaranteed by Assumption~\ref{assm:nondegeneracy}; its gradient is bounded there (scaling like $c_0^{-3/2}$). This makes the non-degeneracy requirement explicit rather than implicit in a ``bounded gradient'' assumption, and localizes where the delta-method linearization degrades (as $c_0 \to 0$).
\end{remark}

\subsection{Post-Proof Audit}

\paragraph{Assumption check.}
All assumptions stated upfront in Section 2.2 and explicitly invoked throughout:
\begin{itemize}
\item Assumptions~\ref{assm:treatment-effects}--\ref{assm:functional}: Treatment effects and functional regularity
\item Assumptions~\ref{assm:nuisance}--\ref{assm:lipschitz}: Nuisance estimation (new for AIPW)
\item Assumptions~\ref{assm:mcmc}--\ref{assm:geometry}: MCMC and geometry
\item Assumption~\ref{assm:sample-size}: Sample size
\end{itemize}

\paragraph{Dependency check.}
All invoked results have their conditions verified:
\begin{itemize}
\item \textbf{Chernozhukov et al.\ (2018) Theorem 3.1:} Conditions verified in Lemma~\ref{lem:cross-fitting} (Neyman orthogonality, $L^2$ rates, cross-fitting, bounded moments)
\item \textbf{CLT:} i.i.d.\ observations, finite variance (Assumption~\ref{assm:bounded-if-variance})
\item \textbf{Delta method:} Hadamard differentiability (Assumption~\ref{assm:functional}), Lipschitz gradient (Assumption~\ref{assm:lipschitz})
\item \textbf{Fubini:} Bounded integrands (Assumptions~\ref{assm:treatment-effects}(c), \ref{assm:functional}, \ref{assm:bounded-if-variance})
\item \textbf{Dominated convergence:} Uniform moment bounds (Assumption~\ref{assm:bounded-if-variance})
\end{itemize}

\paragraph{Rate verification.}
\begin{itemize}
\item \textbf{Nuisance estimation:} $o_P(n^{-1/4})$ for propensity and outcome regressions (Assumption~\ref{assm:nuisance}(b-c))
\item \textbf{Second-order remainder:} $\|\hat{\eta} - \eta_0\|_{L^2}^2 = [o_P(n^{-1/4})]^2 = o_P(n^{-1/2})$ (Assumption~\ref{assm:nuisance}(d))
\item \textbf{Treatment effect estimators:} $O_P(n^{-1/2})$ rate (Lemma~\ref{lem:cross-fitting})
\item \textbf{MCMC approximation:} $O_P(M^{-1/2})$, independent of $n$ (Assumption~\ref{assm:mcmc})
  \begin{itemize}
  \item If $M$ fixed: Non-vanishing bias in target
  \item If $M \to \infty$: Bias vanishes
  \end{itemize}
\item \textbf{Remainder terms:} All $o_P(n^{-1/2})$ by Neyman orthogonality and bounded second derivatives
\item \textbf{Final rate:} $\sqrt{n}$-consistent (conditional on MCMC samples)
\end{itemize}

\paragraph{Key differences from RCT case.}
\begin{itemize}
\item \textbf{Influence function:} Three-term AIPW structure replaces simple weighted difference-in-means
\item \textbf{Nuisance estimation:} Propensity scores and outcome regressions estimated with $o_P(n^{-1/4})$ rates
\item \textbf{Cross-fitting:} Sample splitting prevents overfitting bias
\item \textbf{Neyman orthogonality:} First-order nuisance bias eliminated (Lemma~\ref{lem:orthogonality})
\item \textbf{Double robustness:} Consistency if either $\hat{e}$ or $\hat{\mu}_a$ correct (structural property of AIPW)
\end{itemize}

\paragraph{What stays the same.}
\begin{itemize}
\item MCMC approximation rate: Still $O_P(M^{-1/2})$
\item Functional delta method: Still applies via Hadamard differentiability
\item Final $\sqrt{n}$ rate: Unchanged
\item Bounded weights key to uniformity: Still central to proof
\end{itemize}

\paragraph{Edge case remarks.}
\begin{enumerate}
\item \textbf{Degenerate variance (correlation functional):} If $\text{Var}_\mu(\Delta_S) = 0$ or $\text{Var}_\mu(\Delta_Y) = 0$, the correlation functional is undefined. The proof requires non-degeneracy: $\min(\text{Var}_\mu(\Delta_S), \text{Var}_\mu(\Delta_Y)) \geq c_0 > 0$, stated explicitly as Assumption~\ref{assm:nondegeneracy}.

\item \textbf{Extreme reweighting:} If $\max_i w_i(\mathcal{Q}) \approx C$ (near the bound), variance can inflate significantly. In practice, recommend $C \leq 10$ to avoid extreme variance inflation. Effective sample size drops as $n_{\text{eff}} \approx n / (1 + \text{CV}^2)$ where $\text{CV}$ is coefficient of variation of weights.

\item \textbf{Non-smooth functionals:} Indicator functions (e.g., $\phi(\Delta_S, \Delta_Y) = \mathbb{1}(\Delta_S > 0)$) are not Hadamard differentiable. The proof applies only to smooth functionals: correlation, covariance, products of moments. For non-smooth functionals, bootstrap or other methods required.

\item \textbf{$M \to \infty$ clarification:} Three separate issues with large $M$:
\begin{itemize}
\item \textbf{Remainder uniformity:} No restriction on $M$ growth rate (bound independent of $M$). This is the key innovation.
\item \textbf{CLT for composite IF:} Requires uniform moment bound $\sup_{\mathcal{Q}} \mathbb{E}[\psi^2] < \infty$ (Assumption~\ref{assm:bounded-if-variance}). Without this, CLT can break as $M \to \infty$.
\item \textbf{MCMC approximation quality:} Requires $\|\hat{\mu}_M - \mu\|_{\text{TV}} = O_P(M^{-1/2})$ (Assumption~\ref{assm:mcmc}). About approximation error, not uniformity.
\end{itemize}

\item \textbf{Strong positivity threshold:} Assumption~\ref{assm:nuisance}(a) requires $e(X) \in [0.1, 0.9]$ (stronger than standard $e(X) \in (0, 1)$). This ensures second derivative bound $\frac{2M}{c^3} \leq 2000$ stays reasonable. Near boundaries ($e(X) \approx 0$ or $1$), second derivatives explode, violating Assumption~\ref{assm:second-derivatives}. Use trimming or alternative methods for weak overlap regions.
\end{enumerate}

\paragraph{Tightness assessment.}
The $\sqrt{n}$ rate is sharp (semiparametric efficiency bound for this estimand). The $o_P(n^{-1/4})$ nuisance rates are also sharp under standard smoothness conditions (cannot be weakened without additional structure). Assumptions are near-minimal for these rates.

\paragraph{Failure localization.}
If asymptotic normality fails or coverage is poor:
\begin{enumerate}
\item \textbf{Check propensity overlap:} $\min_i \hat{e}(X_i) \geq 0.1$ and $\max_i \hat{e}(X_i) \leq 0.9$? If not, trim or use alternative methods.
\item \textbf{Check nuisance convergence:} Use cross-validated MSE to verify $o_P(n^{-1/4})$ rates. If slower, coverage will be poor.
\item \textbf{Check weight bounds:} $\max_m \max_i w_i(\mathcal{Q}_m) \leq C$? Large weights inflate variance.
\item \textbf{Check variance degeneracy:} $\min(\text{Var}_\mu(\Delta_S), \text{Var}_\mu(\Delta_Y)) > 0$? If near-zero, correlation undefined.
\item \textbf{Check MCMC convergence:} Is $M$ large enough? Is burn-in sufficient? Use $\hat{R}$ diagnostics.
\item \textbf{Check MCMC bias:} If $M$ is fixed and coverage systematically off, MCMC bias $O_P(M^{-1/2})$ may dominate. Increase $M$.
\end{enumerate}

\paragraph{Implementation notes.}
The R package \texttt{surrogateTransportability} implements this framework:
\begin{itemize}
\item \texttt{tv\_ball\_correlation\_IF\_adaptive()}: main estimator (adaptive hit-and-run,
  importance-weighting or cross-fitted AIPW, influence-function inference)
\item \texttt{sample\_tv\_ball()}: uniform hit-and-run sampling on the TV ball
\item \texttt{gradient\_correlation\_analytical()}: analytic gradient of the correlation functional
\end{itemize}
See package documentation for usage examples.

% =====================================================================
%  GENERAL ASYMPTOTIC-NORMALITY THEORY (extends Theorem \ref{thm:asymptotic-normality})
%  Labels use an "A" scheme (thm:general, lem:*-A, cor:finite-support).
%  Notation matches the finite-support proof above (explicit long form).
% =====================================================================

\section{General Estimand: Asymptotic Normality via the Bilinear-Functional Structure}
\label{sec:general-theory}

This section extends Theorem~\ref{thm:asymptotic-normality} from the finite-support,
correlation-specific case to a general smooth functional $\Theta = \Psi(m_1,\dots,m_K)$
of the across-study effect distribution, on discrete \emph{or} smooth-enough continuous
covariate spaces. As in Theorem~\ref{thm:asymptotic-normality}, we \emph{condition on the
sampled geometry}: the MCMC sample $\hat{\mu}_M$ (equivalently the reweighting-covariance
kernel $C$ it induces, defined below) is treated as fixed, and the $\sqrt{M}$ MCMC term is
handled exactly as before via Lemma~\ref{lem:mcmc-error}.

\subsection{General estimand and the bilinear-functional structure}
\label{subsec:bilinear-structure}

\paragraph{Covariate-level effects.}
Let $O=(X,A,S,Y)\sim\mathbb{P}_0$ with $A\in\{0,1\}$ and $X\in\mathcal{X}$, where
$\mathcal{X}\subseteq\mathbb{R}^d$ (finite, or continuous of dimension $d$). Define the
conditional average treatment effects (CATEs)
\begin{equation}
\tau_S(x) = \mathbb{E}_{\mathbb{P}_0}[S(1)-S(0)\mid X=x],
\qquad
\tau_Y(x) = \mathbb{E}_{\mathbb{P}_0}[Y(1)-Y(0)\mid X=x],
\end{equation}
and write $\beta=(\tau_S,\tau_Y)$. For a future study $\mathcal{Q}\ll\mathbb{P}_0$ with
density ratio $w(x)=\tfrac{d\mathcal{Q}}{d\mathbb{P}_0}(x)$, the study-level effects are the
\emph{linear} functionals
\begin{equation}
\Delta_a(\mathcal{Q}) = \mathbb{E}_{\mathbb{P}_0}[\,w(X)\,\tau_a(X)\,],
\qquad a\in\{S,Y\}.
\label{eq:delta-linear}
\end{equation}

\paragraph{The geometry and its kernel.}
A geometry is a probability measure $\mu$ over reweightings $w$, supported on a local
convex body $\{\mathcal{Q}:d(\mathcal{Q},\mathbb{P}_0)\le\lambda\}$ (Assumption~\ref{assm:geometry}).
Define
\begin{equation}
\bar{a}(x) = \mathbb{E}_\mu[w(x)],
\qquad
C(x,x') = \mathrm{Cov}_\mu\!\big(w(x),\,w(x')\big),
\label{eq:kernel-def}
\end{equation}
the reweighting mean and the \emph{reweighting-covariance kernel}. By construction $C$ is
symmetric; under Assumptions~\ref{assm:treatment-effects}(c) and \ref{assm:geometry} it is
bounded, $\sup_{x,x'}|C(x,x')|\le C_w^2<\infty$.

\paragraph{Linear/bilinear reduction.}
Taking $\mu$-moments of \eqref{eq:delta-linear} and using Fubini (justified by bounded weights
and $\mathbb{E}_{\mathbb{P}_0}[\tau_a^2]<\infty$):
\begin{align}
\mathbb{E}_\mu[\Delta_a(\mathcal{Q})]
   &= \int \bar{a}(x)\,\tau_a(x)\,d\mathbb{P}_0(x),
   \tag{linear}\label{eq:reduction-linear}\\
\mathrm{Cov}_\mu\!\big(\Delta_a(\mathcal{Q}),\Delta_b(\mathcal{Q})\big)
   &= \iint C(x,x')\,\tau_a(x)\,\tau_b(x')\,d\mathbb{P}_0(x)\,d\mathbb{P}_0(x')
   \;=:\; \Theta_{ab},
   \tag{bilinear}\label{eq:reduction-bilinear}
\end{align}
for $a,b\in\{S,Y\}$. Thus every load-bearing $\mu$-moment is either a bounded \emph{linear}
functional of $\beta$ (against the weight $\bar{a}$) or a bounded \emph{bilinear} functional
of $\beta$ (against the kernel $C$). We write the moment vector
\begin{equation}
m = \big(\underbrace{\mathbb{E}_\mu\Delta_S,\ \mathbb{E}_\mu\Delta_Y}_{\text{linear}},\ 
        \underbrace{\Theta_{SS},\ \Theta_{YY},\ \Theta_{SY}}_{\text{bilinear}},\ \dots\big)
        \in\mathbb{R}^K,
\end{equation}
and the estimand is $\Theta = \Psi(m)$ for a smooth $\Psi:\mathbb{R}^K\to\mathbb{R}$
(Assumption~\ref{assm:functional}). Correlation, $R^2$, MSPE, and concordance are instances.

\paragraph{Discrete specialization ($C\to\Sigma$).}
When $X\in\{1,\dots,K\}$ is finite with cell probabilities $p_0(k)$, write
$\tau_a=(\tau_a(1),\dots,\tau_a(K))^\top$ and $q=(q_1,\dots,q_K)^\top$ the future-study cell
mass. Then $w(k)=q_k/p_0(k)$, $\bar{a}(k)=\mathbb{E}_\mu[w(k)]$, and the kernel collapses to a
$K\times K$ matrix $\Sigma$ with entries
$\Sigma_{kk'} = C(k,k')\,p_0(k)\,p_0(k')=\mathrm{Cov}_\mu(q_k,q_{k'})$, so that
\begin{equation}
\Theta_{ab} = \iint C\,\tau_a\tau_b\,d\mathbb{P}_0^2
            = \sum_{k,k'} p_0(k)p_0(k')\,C(k,k')\,\tau_a(k)\tau_b(k')
            = \tau_a^\top \Sigma\, \tau_b .
\label{eq:discrete-quadform}
\end{equation}
This is the $\tau^\top\Sigma\tau$ object of the finite-support proof; Corollary~\ref{cor:finite-support}
shows the general result collapses to Theorem~\ref{thm:asymptotic-normality} here.

\subsection{Assumption ledger A0--A8}
\label{subsec:assumptions-A}

Assumptions \ref{assm:treatment-effects}--\ref{assm:functional} and
\ref{assm:nuisance}--\ref{assm:sample-size} carry over as the general-case A1--A4, A6--A8
(with $\phi$-specific statements reread for $\Psi$ and the moment vector $m$). We state in full
the two assumptions that are new relative to the finite-support proof: A0 (the causal scope of
the estimand) and A5 (the bilinear-functional estimability boundary).

\begin{assumption}[A0: Conditional-effect transportability given $X$]
\label{assm:A0-transport}
All cross-study variation in treatment effects is compositional in the measured covariates
$X$: future studies $\mathcal{Q}$ differ from $\mathbb{P}_0$ only in the marginal law of $X$
and share $\mathbb{P}_0$'s conditional average treatment effects $\tau_S(\cdot),\tau_Y(\cdot)$
given $X$. Equivalently, the reweighting $w=\tfrac{d\mathcal{Q}}{d\mathbb{P}_0}$ is a function
of $X$ alone, so $\tau_S(x),\tau_Y(x)$ are frozen at their $\mathbb{P}_0$ values under every
$\mathcal{Q}$, which is exactly what makes \eqref{eq:delta-linear} hold. (Only the CATEs, not
the full conditional law of potential outcomes, need be shared: \eqref{eq:delta-linear} depends
on $\mathbb{P}_0$ through $(\tau_S,\tau_Y)$ alone.)
\end{assumption}

\begin{remark}[Definitional vs.\ interpretive reading of A0]
\label{rem:A0-reading}
Under the \emph{definitional} reading the estimand simply \emph{is} the functional over the
class of $X$-compositional future studies, and A0 is a scope statement, not an assumption;
under the \emph{interpretive} reading---reading $\Theta$ as ``does the surrogate transport to
real future studies''---A0 is a substantive causal assumption of the effects-transport-given-$X$
class, distinct from the support condition $\mathcal{Q}\ll\mathbb{P}_0$
(Assumption~\ref{assm:treatment-effects}(c)) and not removable for a single-study point estimate.
\end{remark}

\begin{assumption}[A5: Bilinear-functional $\sqrt{n}$-estimability]
\label{assm:A5-elbow}
Each bilinear functional $\Theta_{ab}$ in \eqref{eq:reduction-bilinear} admits a
$\sqrt{n}$-consistent, asymptotically linear (debiased) estimator. Let $\tau_S,\tau_Y$ lie in
H\"older/Sobolev balls of smoothness $s_S,s_Y$ on $\mathcal{X}\subseteq\mathbb{R}^d$. Then, for
the identity/Dirac kernel, $\sqrt{n}$-estimability of $\Theta_{ab}$ holds if and only if the
sum-of-smoothness condition
\begin{equation}
s_S + s_Y > d/2
\label{eq:elbow}
\end{equation}
holds. When \eqref{eq:elbow} fails, the minimax rate is
$n^{-\,4s/(4s+d)}$ (with $s=\tfrac{1}{2}(s_S+s_Y)$ the effective smoothness; this single-$s$
rate expression is the balanced-smoothness reduction of the genuinely two-smoothness minimax
boundary, whereas the elbow \eqref{eq:elbow} itself is exact), strictly slower
than $n^{-1/2}$, and no $\sqrt{n}$-consistent asymptotically normal estimator exists
\cite{BickelRitov1988,BirgeMassart1995,Robins2008HOIF,Robins2017Minimax,Kennedy2024Minimax,McClean2024DCDR}.
\end{assumption}

\begin{remark}[A5 is conservative; finite support is automatic]
\label{rem:A5-conservative}
Condition \eqref{eq:elbow} is the boundary for the \emph{Dirac} kernel. Our kernel $C$ is
bounded and (typically) smooth under Assumptions~\ref{assm:treatment-effects}(c) and
\ref{assm:geometry}; a genuinely smooth $C$ typically \emph{absorbs} smoothness (its integral
operator regularizes $h_b=\int C\tau_b\,d\mathbb{P}_0$) and thus \emph{relaxes} the requirement,
so \eqref{eq:elbow} is, at worst, a conservative sufficient condition here. On finite
$X$ every function is arbitrarily smooth and $\Theta_{ab}=\tau_a^\top\Sigma\tau_b$ is a fixed
smooth (indeed polynomial) function of finitely many cell means, so A5 holds automatically---%
this is why the finite-support proof never states it. Higher-order influence functions attain
the minimax rate below the elbow but do \emph{not} restore $\sqrt{n}$ normality.
\end{remark}

\subsection{Layer L1: CATE asymptotic linearity}

\begin{lemma}[A1: CATE asymptotic linearity]
\label{lem:cate-linearity-A}
Under A1--A2 (Assumption~\ref{assm:treatment-effects} generalized off the finite simplex, plus
a pluggable estimator supplying a CATE influence function), for each $a\in\{S,Y\}$ and each
$x\in\mathcal{X}$ the CATE estimator admits the expansion
\begin{equation}
\hat{\tau}_a(x) - \tau_a(x)
  = \frac{1}{n}\sum_{i=1}^n \mathrm{IF}_{\tau_a}(O_i;x) + R_a(x),
\label{eq:cate-expansion}
\end{equation}
where $\mathbb{E}_{\mathbb{P}_0}[\mathrm{IF}_{\tau_a}(O;x)]=0$ for every $x$, and the remainder
satisfies $\|R_a\|_{L^2(\mathbb{P}_0)}=o_P(n^{-1/4})$ (product-rate control under A6).
For an RCT $\mathrm{IF}_{\tau_a}$ is the (per-$x$) difference-of-means score; for AIPW /
DR-learners it is the pointwise AIPW score, and a generic learner is admissible whenever
it supplies a valid $\mathrm{IF}_{\tau_a}$ satisfying \eqref{eq:cate-expansion}.
\end{lemma}

\begin{proof}[Proof of Lemma~\ref{lem:cate-linearity-A}]
This is the pointwise (per-$x$) restatement of the AIPW expansion established, integrated
against $w$, in Lemma~\ref{lem:cross-fitting}; the reweighted per-study score
$\psi_a^{\text{AIPW}}(O;\eta_0,\mathcal{Q})$ of \eqref{eq:aipw-s}--\eqref{eq:aipw-y} equals
$\int w(x)\,\mathrm{IF}_{\tau_a}(O;x)\,d\mathbb{P}_0(x)$ up to the mean-zero centering, so
mean-zero-ness and the $o_P(n^{-1/2})$ integrated remainder transfer directly. For the discrete
case the score is exactly the per-cell AIPW score of the companion derivation
(\texttt{derivation\_influence\_functions.md}, \S2). The remainder rate is A6.
% NOTE(auditor): \eqref{eq:cate-expansion} is stated pointwise in x. Downstream (Lemma
% \ref{lem:bilinear-eif-A}) we only ever integrate IF against dP0 or the bounded kernel C, so a
% pointwise-in-x, L2-in-O influence function suffices; no Donsker/uniform-in-x class is claimed.
\end{proof}

\subsection{Layer L3: the bilinear-functional EIF (core new result)}

\begin{lemma}[A3: EIF of a bilinear functional of CATEs]
\label{lem:bilinear-eif-A}
Fix a bounded symmetric kernel $C$ and the covariate law of $\mathbb{P}_0$; treat both as
known (conditioning on the sampled geometry, as in Theorem~\ref{thm:asymptotic-normality}).
Suppose $\tau_a,\tau_b$ are pathwise differentiable with influence functions
$\mathrm{IF}_{\tau_a},\mathrm{IF}_{\tau_b}$ (Lemma~\ref{lem:cate-linearity-A}), and
$\mathbb{E}_{\mathbb{P}_0}[\tau_a^2],\mathbb{E}_{\mathbb{P}_0}[\tau_b^2]<\infty$. Then
$\Theta_{ab}=\iint C(x,x')\tau_a(x)\tau_b(x')\,d\mathbb{P}_0(x)d\mathbb{P}_0(x')$ is pathwise
differentiable with (mean-zero) efficient influence function
\begin{equation}
\chi_{ab}(O)
 = \int h_b(x)\,\mathrm{IF}_{\tau_a}(O;x)\,d\mathbb{P}_0(x)
 + \int h_a(x')\,\mathrm{IF}_{\tau_b}(O;x')\,d\mathbb{P}_0(x'),
\label{eq:chi-ab}
\end{equation}
where
\begin{equation}
h_b(x) = \int C(x,x')\,\tau_b(x')\,d\mathbb{P}_0(x'),
\qquad
h_a(x') = \int C(x',x)\,\tau_a(x)\,d\mathbb{P}_0(x).
\label{eq:h-directions}
\end{equation}
Since $\mathbb{E}_{\mathbb{P}_0}[\mathrm{IF}_{\tau_a}(O;x)]=0$ for every $x$,
$\mathbb{E}_{\mathbb{P}_0}[\chi_{ab}]=0$; the plug-in value $\Theta_{ab}$ is carried by the
plug-in term of the one-step estimator \eqref{eq:onestep}, not by the influence function.
In particular, when $a=b$,
$\chi_{aa}(O) = 2\int h_a(x)\,\mathrm{IF}_{\tau_a}(O;x)\,d\mathbb{P}_0(x)$ with
$h_a(x)=\int C(x,x')\tau_a(x')d\mathbb{P}_0(x')$; and for finite $X$, \eqref{eq:chi-ab}
reduces to a quadratic/trace form in $\Sigma$ (Corollary~\ref{cor:finite-support}).
\end{lemma}

\begin{proof}[Proof of Lemma~\ref{lem:bilinear-eif-A}]
\emph{Strategy.} $\Theta_{ab}$ depends on $\mathbb{P}_0$ (with $C$ and the covariate law
fixed) only through the pair $(\tau_a,\tau_b)$. We compute the pathwise derivative by the
product rule over the two $\tau$-factors, identify the Riesz representer of each directional
derivative, and symmetrize.

\emph{Step 1: one-dimensional submodel.} Let $\{\mathbb{P}_t\}$ be a regular parametric submodel
through $\mathbb{P}_0$ at $t=0$ with score $g(O)$, and let $\tau_{a,t},\tau_{b,t}$ denote the
CATEs under $\mathbb{P}_t$. Pathwise differentiability of the CATEs (Lemma~\ref{lem:cate-linearity-A})
gives, for each fixed $x$,
\begin{equation}
\frac{d}{dt}\tau_{a,t}(x)\Big|_{t=0}
   = \mathbb{E}_{\mathbb{P}_0}\!\big[\mathrm{IF}_{\tau_a}(O;x)\,g(O)\big],
\label{eq:tau-derivative}
\end{equation}
and similarly for $\tau_b$. (Here $C$ and $d\mathbb{P}_0$ in the definition of $\Theta_{ab}$ are
held fixed; see Remark~\ref{rem:kernel-dependence}.)

\emph{Step 2: product rule over the two factors.} Write
$\Theta_{ab}(t)=\iint C(x,x')\tau_{a,t}(x)\tau_{b,t}(x')\,d\mathbb{P}_0(x)d\mathbb{P}_0(x')$
(covariate law fixed). Differentiating under the integral (Leibniz; justified because $C$ is
bounded, $\tau_{a,t},\tau_{b,t}$ have square-integrable derivatives, and the domination is by
$C_w^2(\,|\partial_t\tau_a||\tau_b| + |\tau_a||\partial_t\tau_b|)\in L^1$):
\begin{align}
\frac{d}{dt}\Theta_{ab}(t)\Big|_{t=0}
 &= \iint C(x,x')\,\Big[\dot\tau_a(x)\,\tau_b(x') + \tau_a(x)\,\dot\tau_b(x')\Big]
     \,d\mathbb{P}_0(x)d\mathbb{P}_0(x') \nonumber\\
 &= \int \dot\tau_a(x)\,h_b(x)\,d\mathbb{P}_0(x)
  + \int \dot\tau_b(x')\,h_a(x')\,d\mathbb{P}_0(x'),
\label{eq:productrule}
\end{align}
where $\dot\tau_a(x)=\tfrac{d}{dt}\tau_{a,t}(x)|_{t=0}$ and we used Fubini (bounded $C$,
square-integrable $\tau$) to collapse each double integral, defining the two \emph{directional}
representers $h_b,h_a$ of \eqref{eq:h-directions}. This is the ``two directional derivatives,
then symmetrize'' step: the $a$-factor is differentiated against the $b$-weight $h_b$, and the
$b$-factor against the $a$-weight $h_a$.

\emph{Step 3: substitute the CATE scores.} Insert \eqref{eq:tau-derivative} into
\eqref{eq:productrule} and interchange the (finite-measure, bounded-integrand) $x$- and
$O$-integrals by Fubini:
\begin{align}
\frac{d}{dt}\Theta_{ab}(t)\Big|_{t=0}
 &= \mathbb{E}_{\mathbb{P}_0}\!\Big[\Big(\int h_b(x)\mathrm{IF}_{\tau_a}(O;x)d\mathbb{P}_0(x)
     + \int h_a(x')\mathrm{IF}_{\tau_b}(O;x')d\mathbb{P}_0(x')\Big)\,g(O)\Big].
\label{eq:riesz}
\end{align}
\emph{Step 4: identify the EIF.} Equation \eqref{eq:riesz} exhibits the pathwise derivative as
$\mathbb{E}_{\mathbb{P}_0}[\chi_{ab}(O)g(O)]$ for every score $g$, where $\chi_{ab}$ is the
bracketed term of \eqref{eq:chi-ab}; by the Riesz representation of pathwise derivatives
$\chi_{ab}$ is the (canonical-gradient) influence function of $\Theta_{ab}$. Since
$\mathbb{E}_{\mathbb{P}_0}[\mathrm{IF}_{\tau_a}(O;x)]=0$ for all $x$, $\chi_{ab}$ is mean-zero,
as an influence function must be; the plug-in value $\Theta_{ab}$ enters not $\chi_{ab}$ but the
plug-in term of the one-step estimator \eqref{eq:onestep}. Efficiency holds because, in the
nonparametric model where the tangent space is all of $L^2_0(\mathbb{P}_0)$, $\chi_{ab}$ is a
$\mathbb{P}_0$-integral of the efficient CATE influence functions and hence lies in the tangent
space, so it is the canonical gradient.

\emph{Step 5: diagonal check.} If $a=b$ then $h_a=h_b$ and $\mathrm{IF}_{\tau_a}=\mathrm{IF}_{\tau_b}$,
so \eqref{eq:chi-ab} becomes
$\chi_{aa}(O)=2\int h_a(x)\mathrm{IF}_{\tau_a}(O;x)d\mathbb{P}_0(x)$, the factor $2$
being exactly the two identical directional derivatives---consistent with $\Theta_{aa}$ being a
pure quadratic whose derivative carries the classical factor $2$.

\emph{Step 6: discrete check.} On finite $X$, \eqref{eq:h-directions} gives
$h_b(k)=\sum_{k'}C(k,k')\tau_b(k')p_0(k')$; substituting into \eqref{eq:chi-ab} and using
$\int h_b\,\mathrm{IF}_{\tau_a}d\mathbb{P}_0=\sum_k p_0(k)h_b(k)\mathrm{IF}_{\tau_a}(\cdot;k)$
reproduces the discrete $\Sigma$-quadratic score; contracting against the CATE-score covariance
$V$ yields the $\mathrm{tr}(\Sigma V)$ debiasing term of the finite-support case
(Corollary~\ref{cor:finite-support}).
\end{proof}

\begin{remark}[Scope: kernel and covariate-law dependence held fixed]
\label{rem:kernel-dependence}
Lemma~\ref{lem:bilinear-eif-A} differentiates only through $(\tau_a,\tau_b)$, holding the
kernel $C$ and the covariate law $d\mathbb{P}_0$ fixed. This matches
Theorem~\ref{thm:asymptotic-normality}, which conditions on the sampled geometry $\hat{\mu}_M$
(hence on $C$). The \emph{unconditional} EIF, which would additionally differentiate $C$'s and
the covariate law's dependence on $\mathbb{P}_0$, contributes further terms; deriving it is
future work and is not needed for the conditional statement of Theorem~\ref{thm:general}.
% NOTE(auditor): Genuine deferral, not a hidden step. Conditional-on-geometry stance inherited
% from Theorem \ref{thm:asymptotic-normality}. An unconditional statement would add Gateaux terms
% from differentiating dP0(x), dP0(x'), and (if C is a P0-plug-in) C itself. Localize here.
\end{remark}

\subsection{Layer L3: one-step debiased cross-fit estimator}

\begin{lemma}[A4: debiased cross-fit estimator of $\Theta_{ab}$]
\label{lem:onestep-A}
Let $\hat\tau_a,\hat\tau_b,\hat{C}$ be cross-fitted estimators and
$\hat h_b(x)=\int \hat C(x,x')\hat\tau_b(x')d\mathbb{P}_0(x')$ (etc.). The one-step estimator
\begin{equation}
\hat\Theta_{ab}
 = \iint \hat C(x,x')\hat\tau_a(x)\hat\tau_b(x')\,d\mathbb{P}_0(x)d\mathbb{P}_0(x')
 + \frac{1}{n}\sum_{i=1}^n \hat\chi_{ab}(O_i)
\label{eq:onestep}
\end{equation}
is Neyman-orthogonal in $\beta=(\tau_a,\tau_b)$: the Gateaux derivative of the map
$\beta\mapsto \mathbb{E}_{\mathbb{P}_0}[\text{(plug-in)}+\chi_{ab}]$ vanishes at the truth.
\end{lemma}

\begin{proof}[Proof of Lemma~\ref{lem:onestep-A}]
By Lemma~\ref{lem:bilinear-eif-A}, $\chi_{ab}$ is the canonical gradient of $\Theta_{ab}$ in
$\beta$; adding it to the plug-in cancels the first-order plug-in bias, the defining property of
the one-step construction. Concretely, the directional derivative of the plug-in at $\beta$
along $(\delta_a,\delta_b)$ is $\int h_b\delta_a\,d\mathbb{P}_0+\int h_a\delta_b\,d\mathbb{P}_0$
by \eqref{eq:productrule}, while $\mathbb{E}_{\mathbb{P}_0}[\chi_{ab}]$ has directional derivative
equal to its negative; the two cancel. The representer $h_b$ plays the role of the weight $w$ in
Lemma~\ref{lem:orthogonality}: since $h_b,h_a$ are functions of $x$ that do \emph{not} depend on
the outcome-model nuisances $\eta=(e,\mu_0,\mu_1)$ (only on $C$ and $\tau$), the
``orthogonality survives reweighting'' argument of Lemma~\ref{lem:orthogonality}, Step~4, applies
verbatim with $h_b$ (resp.\ $h_a$) in the role of $w(X;\mathcal{Q})$, so orthogonality \emph{in
$\eta$} holds in addition to orthogonality \emph{in $\beta$}. (The estimated $\hat h_b=\int\hat
C\hat\tau_b\,d\mathbb{P}_0$ depends on $\eta$ only through $\hat\tau_b$, which is the
$\beta$-argument handled by $\beta$-orthogonality, not through a separate $\eta$-path, since
$\hat C$ is geometry and hence $\eta$-free; the inheritance is therefore airtight.)
% NOTE(auditor): Two orthogonalities: (i) in beta by the one-step construction; (ii) in eta,
% inherited from Lemma \ref{lem:orthogonality} because h_a,h_b are eta-free x-functions. Both are
% needed for the Lemma \ref{lem:elbow} remainder to be second order. Most worth re-checking.
\end{proof}

\subsection{Layer L3: remainder and the estimability elbow}

\begin{lemma}[A5: bilinear remainder / elbow]
\label{lem:elbow}
Under Lemmas~\ref{lem:cate-linearity-A}, \ref{lem:onestep-A}, the estimator
\eqref{eq:onestep} satisfies
$\hat\Theta_{ab} - \Theta_{ab} = \frac{1}{n}\sum_{i=1}^n \chi_{ab}(O_i) + R_n$,
where the second-order remainder is, up to $o_P(n^{-1/2})$ nuisance terms controlled by A6 and
Lemma~\ref{lem:onestep-A},
\begin{equation}
R_n = \iint \hat C(x,x')\,(\hat\tau_a-\tau_a)(x)\,(\hat\tau_b-\tau_b)(x')\,d\mathbb{P}_0(x)d\mathbb{P}_0(x').
\label{eq:remainder}
\end{equation}
By Cauchy--Schwarz with the bounded symmetric kernel $\hat C$ (operator norm
$\|\hat C\|_{op}\le C_w^2$ under Assumptions~\ref{assm:treatment-effects}(c), \ref{assm:geometry}),
\begin{equation}
|R_n|
 \le \|\hat C\|_{op}\,\big\|\hat\tau_a-\tau_a\big\|_{L^2(\mathbb{P}_0)}\,
     \big\|\hat\tau_b-\tau_b\big\|_{L^2(\mathbb{P}_0)}.
\label{eq:CS-remainder}
\end{equation}
Hence $R_n=o_P(n^{-1/2})$ if $\|\hat\tau_a-\tau_a\|\cdot\|\hat\tau_b-\tau_b\| = o_P(n^{-1/2})$,
which holds under A5: for H\"older/Sobolev smoothness $s_S,s_Y$ the achievable rates give a
product $o_P(n^{-1/2})$ exactly when $s_S+s_Y>d/2$. When \eqref{eq:elbow} fails, $\hat\Theta_{ab}$
has the slower minimax rate $n^{-4s/(4s+d)}$ and no $\sqrt{n}$-normal estimator exists
\cite{BickelRitov1988,BirgeMassart1995,Robins2008HOIF,Kennedy2024Minimax,McClean2024DCDR}.
\end{lemma}

\begin{proof}[Proof of Lemma~\ref{lem:elbow}]
The von Mises expansion of \eqref{eq:onestep} has exact second-order remainder
\eqref{eq:remainder}: the one-step correction removes the first-order term
(Lemma~\ref{lem:onestep-A}), and $\Theta_{ab}$ is exactly \emph{bilinear} in $\beta$, so the
Taylor expansion terminates at the product term with \emph{no} higher-order terms (a quadratic
functional has zero third derivative). The nuisance-$\eta$ contribution is $o_P(n^{-1/2})$ by
Neyman orthogonality in $\eta$ (Lemma~\ref{lem:onestep-A}) and the product-rate A6, exactly as in
Lemma~\ref{lem:cross-fitting}. For \eqref{eq:CS-remainder}: writing $\delta_a=\hat\tau_a-\tau_a$,
$\delta_b=\hat\tau_b-\tau_b$, $R_n=\langle \delta_a,\hat C\delta_b\rangle_{L^2(\mathbb{P}_0)}$
with $(\hat C\delta_b)(x)=\int\hat C(x,x')\delta_b(x')d\mathbb{P}_0(x')$; Cauchy--Schwarz and the
operator-norm bound give $|R_n|\le\|\hat C\|_{op}\|\delta_a\|\|\delta_b\|$. The smoothness
accounting for the product rate is the quadratic-functional elbow of A5.
\end{proof}

\subsection{Layers L3--L4: moment-vector EIF and composition}

\begin{lemma}[A7: moment-vector EIF and delta-method composition]
\label{lem:moment-eif-A}
Under A0--A8, the moment vector $m$ has influence function $\psi_m$ obtained by stacking:
\begin{enumerate}[(a)]
\item the linear-moment scores
$\psi_{\mathbb{E}_\mu\Delta_a}(O)=\int\bar a(x)\,\mathrm{IF}_{\tau_a}(O;x)\,d\mathbb{P}_0(x)$
for $a\in\{S,Y\}$;
\item the bilinear scores $\chi_{SS},\chi_{YY},\chi_{SY}$ from Lemma~\ref{lem:bilinear-eif-A};
\item for raw second moments, using $\mathbb{E}_\mu[\Delta_a^2]=\Theta_{aa}+(\mathbb{E}_\mu\Delta_a)^2$, the EIF
$\psi_{\mathbb{E}_\mu\Delta_a^2}=\chi_{aa}+2\,\mathbb{E}_\mu[\Delta_a]\,\psi_{\mathbb{E}_\mu\Delta_a}$.
\end{enumerate}
Then, for smooth $\Psi$ with gradient $\nabla\Psi(m)$ (A7 non-degeneracy),
\begin{equation}
\psi_\Theta(O) = \nabla\Psi(m)^\top\,\psi_m(O).
\label{eq:psi-theta}
\end{equation}
The cross-outcome score $\chi_{SY}$ carries
$\mathbb{E}_{\mathbb{P}_0}[\mathrm{IF}_{\tau_S}(O;\cdot)\,\mathrm{IF}_{\tau_Y}(O;\cdot)]$
(the two CATE influence functions are evaluated on the \emph{same unit} $O$), so
$\mathrm{Var}(\psi_\Theta)$ is \textbf{not} block-diagonal across $S$ and $Y$.
\end{lemma}

\begin{proof}[Proof of Lemma~\ref{lem:moment-eif-A}]
(a) is the linear-functional EIF (mean-zero and Riesz representation as in
Lemma~\ref{lem:bilinear-eif-A}, Step~4, with the degenerate weight $\bar a$). (b) is
Lemma~\ref{lem:bilinear-eif-A}. (c) is the chain rule on
$m\mapsto\Theta_{aa}+(\mathbb{E}_\mu\Delta_a)^2$. Finally \eqref{eq:psi-theta} is the functional
delta method (Lemma~\ref{lem:functional-expansion}) applied to $\Psi$ at $m$, whose Hadamard
differentiability and bounded gradient are A7 (Assumptions~\ref{assm:functional},
\ref{assm:lipschitz}, \ref{assm:nondegeneracy}). Non-block-diagonality is immediate because
$\chi_{SY}$ pairs $\mathrm{IF}_{\tau_S}$ and $\mathrm{IF}_{\tau_Y}$ at the same $O$.
\end{proof}

\subsection{Main result: general asymptotic normality}

\begin{theorem}[A: general asymptotic normality, conditional on the geometry]
\label{thm:general}
Under Assumptions~\ref{assm:A0-transport} (A0), \ref{assm:A5-elbow} (A5), and the general-case
readings of Assumptions~\ref{assm:treatment-effects}--\ref{assm:functional} (A1--A4),
\ref{assm:nuisance}--\ref{assm:sample-size} (A6--A8), with one-step cross-fit debiased
estimators of the moments $m$ (Lemmas~\ref{lem:cate-linearity-A}--\ref{lem:elbow}) and
$\hat\Theta=\Psi(\hat m)$, conditional on the sampled geometry $\hat{\mu}_M$ (equivalently on
the kernel $C$), as $n\to\infty$,
\begin{equation}
\sqrt{n}\big(\hat\Theta - \Theta(\hat{\mu}_M)\big)
 = \frac{1}{\sqrt{n}}\sum_{i=1}^n \psi_\Theta(O_i;\hat{\mu}_M)
 + O_P\!\big((n/M)^{1/2}\big) + o_P(1)
 \xrightarrow{d} N\!\big(0,\ \mathbb{E}_{\mathbb{P}_0}[\psi_\Theta^2]\big),
\end{equation}
with $\psi_\Theta$ of \eqref{eq:psi-theta}. If additionally $n=o(M)$
(Assumption~\ref{assm:sample-size}), the MCMC term is $o_P(1)$ and unconditionally
$\sqrt{n}(\hat\Theta-\Theta(\mu))\xrightarrow{d}N(0,\sigma^2(\mu))$,
$\sigma^2(\mu)=\mathbb{E}_{\mathbb{P}_0}[\psi_\Theta(O;\mu)^2]$.
\end{theorem}

\begin{proof}[Proof of Theorem~\ref{thm:general}]
\emph{Step 1 (moment expansions).} For the linear moments, Lemma~\ref{lem:cate-linearity-A}
integrated against $\bar a$ gives
$\widehat{\mathbb{E}_\mu\Delta_a}-\mathbb{E}_\mu\Delta_a=\frac1n\sum_i\psi_{\mathbb{E}_\mu\Delta_a}(O_i)+o_P(n^{-1/2})$.
For the bilinear moments, Lemma~\ref{lem:elbow} gives
$\hat\Theta_{ab}-\Theta_{ab}=\frac1n\sum_i\chi_{ab}(O_i)+o_P(n^{-1/2})$ \emph{provided A5 holds}
(the only use of A5; below the elbow this step fails and the rate degrades). Stacking,
$\hat m-m=\frac1n\sum_i\psi_m(O_i)+o_P(n^{-1/2})$.

\emph{Step 2 (delta method / L4).} $\Psi$ is Hadamard differentiable with bounded gradient on
$\{\mathrm{Var}_\mu\Delta_S,\mathrm{Var}_\mu\Delta_Y\ge c_0\}$ (Assumption~\ref{assm:nondegeneracy});
Lemma~\ref{lem:functional-expansion} (reused, with $m$ in place of the per-study pair) gives
$\sqrt n(\hat\Theta-\Theta(\hat\mu_M)) = \nabla\Psi(m)^\top\sqrt n(\hat m-m)+o_P(1)
= \frac1{\sqrt n}\sum_i\psi_\Theta(O_i;\hat\mu_M)+o_P(1)$ by Lemma~\ref{lem:moment-eif-A}. For the
ratio-type $\Psi$ (correlation), the second-order delta-method remainder is
$O_P(\|\hat m-m\|^2)=O_P(n^{-1})=o_P(n^{-1/2})$ since $\nabla^2\Psi$ is bounded on the
non-degenerate region (A7).
% NOTE(auditor): ratio-Psi second-order remainder closes iff bounded Hessian x O_P(n^{-1}) =
% o_P(n^{-1/2}), i.e. under A7 non-degeneracy. If c_0 -> 0 a TMLE targeting step substitutes for
% one-step + delta method. This is the ratio-sufficiency point to re-check.

\emph{Step 3 (CLT, conditional on $\hat\mu_M$).} Conditional on the geometry, $O_1,\dots,O_n$ are
i.i.d.\ and independent of $\hat\mu_M$; $\mathbb{E}_{\mathbb{P}_0}[\psi_\Theta]=0$ (each stacked
score is mean-zero: Lemma~\ref{lem:bilinear-eif-A}, Step~4) and
$\mathbb{E}_{\mathbb{P}_0}[\psi_\Theta^2]<\infty$ under Assumption~\ref{assm:bounded-if-variance}
(bounded weights, bounded $C$, bounded gradient). The Lindeberg--L\'evy CLT gives
$\frac1{\sqrt n}\sum_i\psi_\Theta(O_i;\hat\mu_M)\xrightarrow{d}N(0,\mathbb{E}[\psi_\Theta^2])$.

\emph{Step 4 (MCMC term).} $\Theta(\hat\mu_M)-\Theta(\mu)=O_P(M^{-1/2})$ by
Lemma~\ref{lem:mcmc-error} applied to the bounded functional $g=\Psi(m(\cdot))$ of the chain
(bounded under A3--A4, A7); scaling by $\sqrt n$ gives the $O_P((n/M)^{1/2})$ term, which is
$o_P(1)$ iff $n=o(M)$ (Assumption~\ref{assm:sample-size}), yielding the unconditional statement.
Combining Steps 1--4 gives the display.
\end{proof}

\subsection{Finite-support special case}

\begin{corollary}[A10: recovery of Theorem~\ref{thm:asymptotic-normality}]
\label{cor:finite-support}
When $X\in\{1,\dots,K\}$ is finite, $C\to\Sigma$ (with $\Sigma_{kk'}=\mathrm{Cov}_\mu(q_k,q_{k'})$),
$\mathrm{IF}_{\tau_a}$ is the per-cell AIPW score, and $\Psi$ is the correlation map $\varphi$ of
Section~\ref{sec:correlation-application}, Theorem~\ref{thm:general} reduces to
Theorem~\ref{thm:asymptotic-normality}, with the same influence function $\psi_\Theta$ and the
same $\sigma^2(\hat\mu_M)$; the $\chi_{ab}$ centering becomes the discrete $\mathrm{tr}(\Sigma V)$
debiasing form, where $V_{kk'}=\mathbb{E}_{\mathbb{P}_0}[\mathrm{IF}_{\tau_a}(O;k)\mathrm{IF}_{\tau_b}(O;k')]$
is the CATE-score covariance.
\end{corollary}

\begin{proof}[Proof of Corollary~\ref{cor:finite-support}]
On finite $X$, A5 is automatic (Remark~\ref{rem:A5-conservative}) and
$\Theta_{ab}=\tau_a^\top\Sigma\tau_b$ by \eqref{eq:discrete-quadform}. The directional representer
satisfies $(h_b(k)p_0(k))_k \propto \Sigma\tau_b$, so Lemma~\ref{lem:bilinear-eif-A} gives the
(mean-zero) score
$\chi_{ab}(O)=(\Sigma\tau_b)^\top\mathbf{IF}_{\tau_a}(O)+(\Sigma\tau_a)^\top\mathbf{IF}_{\tau_b}(O)$,
with $\mathbf{IF}_{\tau_a}(O)=(\mathrm{IF}_{\tau_a}(O;1),\dots,\mathrm{IF}_{\tau_a}(O;K))^\top$
(the plug-in $\tau_a^\top\Sigma\tau_b$ is carried by the plug-in term of \eqref{eq:onestep}).
The population mean of the first-order plug-in bias contracts the CATE-score vectors against
$\Sigma$, giving $\mathrm{tr}(\Sigma V)$ with
$V=\mathbb{E}_{\mathbb{P}_0}[\mathbf{IF}_{\tau_a}\mathbf{IF}_{\tau_b}^\top]$---the disattenuation
term of the finite-support estimator. The per-cell AIPW score is that of
Lemma~\ref{lem:cross-fitting} restricted to a cell, and $\varphi$'s gradient is the composed
gradient of Remark~\ref{rem:aggregate-gradient}; substituting into \eqref{eq:psi-theta} reproduces
$\psi_\Theta(O;\hat\mu_M)$ of Theorem~\ref{thm:asymptotic-normality}. All remainders are those of
Lemmas~\ref{lem:cross-fitting}, \ref{lem:functional-expansion}, so the two limit laws coincide.
% NOTE(auditor): The exact tr(Sigma V) constant depends on the factor-2 diagonal vs off-diagonal
% pairing (Lemma \ref{lem:bilinear-eif-A}, Step 5). The claim is STRUCTURAL collapse (same object),
% not a re-derivation of the finite-support constant (already validated in
% Theorem \ref{thm:asymptotic-normality}). Re-check pairing if a numeric match is asserted.
\end{proof}

\subsection{Post-proof audit (general theory)}
\label{subsec:audit-general}

\paragraph{Assumption check.}
A0 (Assumption~\ref{assm:A0-transport}) is used once, to license the frozen-effect
representation \eqref{eq:delta-linear}. A1--A2 enter Lemma~\ref{lem:cate-linearity-A}; A3
(Assumption~\ref{assm:treatment-effects}(c)) bounds $C$ and the weights; A4
(Assumption~\ref{assm:geometry}) gives the bounded kernel and MCMC-CLT; A5
(Assumption~\ref{assm:A5-elbow}) is used once, in Lemma~\ref{lem:elbow}; A6 controls nuisance
remainders; A7 (Assumptions~\ref{assm:functional}, \ref{assm:lipschitz}, \ref{assm:nondegeneracy})
gives the delta method and bounded Hessian; A8 (Assumption~\ref{assm:sample-size}) governs the $n$
vs $M$ regime. No assumption is introduced mid-proof.

\paragraph{Dependency check.}
Riesz representation (Lemma~\ref{lem:bilinear-eif-A}): square-integrable CATE scores and bounded
$C$, verified. Fubini/Leibniz: bounded $C$, $L^2$ CATEs, verified. Cauchy--Schwarz with
operator-norm bound (Lemma~\ref{lem:elbow}): symmetric bounded kernel, verified under A3--A4.
Chernozhukov et al.\ (2018) orthogonality/cross-fitting: reused via Lemmas~\ref{lem:orthogonality},
\ref{lem:cross-fitting} with the $\eta$-free $h_a,h_b$ in the role of $w$. Markov-chain CLT
(Lemma~\ref{lem:mcmc-error}) and delta method (Lemma~\ref{lem:functional-expansion}) reused
unchanged. Minimax lower bounds
\cite{BickelRitov1988,BirgeMassart1995,Robins2008HOIF,Robins2017Minimax,Kennedy2024Minimax,McClean2024DCDR}
supply the below-elbow non-existence claim.

\paragraph{Rate verification.}
Bilinear remainder \eqref{eq:CS-remainder}: product of CATE $L^2$-rates, $o_P(n^{-1/2})$ iff
$s_S+s_Y>d/2$ (A5); at equality $O_P(n^{-1/2})$ (no slack), below strictly slower and dominant, so
no $\sqrt n$-normal limit. Nuisance-$\eta$ remainder: $o_P(n^{-1/2})$ via A6, as in
Lemma~\ref{lem:cross-fitting}. Delta-method second-order remainder: $O_P(n^{-1})=o_P(n^{-1/2})$
under bounded Hessian (A7). MCMC: $O_P((n/M)^{1/2})$, $o_P(1)$ iff $n=o(M)$.

\paragraph{Edge-case scan.}
(i) \emph{Near-elbow} $s_S+s_Y\downarrow d/2$: product rate loses slack, intervals widen; report
the minimax rate honestly. (ii) \emph{Weak overlap} $e(x)\to0,1$: $\mathrm{IF}_{\tau_a}$ blows up
(Assumption~\ref{assm:second-derivatives} fails); trim. (iii) \emph{Degenerate variance}: $\nabla\Psi$
blows up like $c_0^{-3/2}$; A7 excludes. (iv) \emph{Kernel conditioning}: the bound
\eqref{eq:CS-remainder} uses $\|C\|_{op}$, not $C^{-1}$, so ill-conditioning of $C$ does not break
it---only extreme weights (bounded by A3) matter.

\paragraph{Tightness.}
Above the elbow the $\sqrt n$ rate is sharp: $\psi_\Theta$ is the canonical gradient, so
$\mathbb{E}[\psi_\Theta^2]$ is the semiparametric efficiency bound. Below the elbow the
$n^{-4s/(4s+d)}$ rate is minimax, so A5 is not improvable as a $\sqrt n$-condition for the Dirac
kernel; for the smooth kernel $C$ it is conservative (Remark~\ref{rem:A5-conservative}).

\paragraph{Failure localization.}
Three steps warrant re-checking. (1) The $a\ne b$ symmetrization
(Lemma~\ref{lem:bilinear-eif-A}, Steps~2--4): the two directional representers must be added,
guarded by the diagonal (Step~5) and discrete (Step~6) checks. (2) Kernel/covariate-law
$\mathbb{P}_0$-dependence (Remark~\ref{rem:kernel-dependence}): explicitly deferred. (3) Ratio
one-step sufficiency (Theorem~\ref{thm:general}, Step~2): closes under bounded Hessian (A7); TMLE
fallback if $c_0\to0$. The linear/bilinear/CLT core is otherwise standard.
